\documentclass[%
  a4paper,%
  11pt,%
  twoside=false,%
  rgb,%
  extramargin,%
  parskip=half,%
]{tubsbook}

\usepackage[utf8x]{inputenc}
\usepackage{xspace}
% \usepackage[T1]{fontenc}
% \usepackage{ae}
% \usepackage[a4paper]{tubsdoc}

\usepackage{scrpage2}
\usepackage{amsmath}
\usepackage{tabularx}
\usepackage{booktabs}
\usepackage{listings}
\lstset{basicstyle=\ttfamily}
\usepackage[%
  colorlinks=true,
  linkcolor=tuRed100,
  citecolor=tuGreenDark]{hyperref}
\usepackage[ngerman]{babel}
\usepackage{caption}
\usepackage{subcaption}
\usepackage{longtable}

% glossaries
\usepackage[ngerman]{translator}
%Paket laden
\usepackage[
nonumberlist, %keine Seitenzahlen anzeigen
acronym,      %ein Abkürzungsverzeichnis erstellen
toc,          %Einträge im Inhaltsverzeichnis
section]      %im Inhaltsverzeichnis auf section-Ebene erscheinen
{glossaries}
%Glossar-Befehle anschalten
\makeglossaries

% Kopiert von KOMA
\makeatletter
\providecommand\marg[1]{%
  {\ttfamily\char`\{}\meta{#1}{\ttfamily\char`\}}}
\providecommand\oarg[1]{%
  {\ttfamily[}\meta{#1}{\ttfamily]}}
\def\cmd#1{\cs{\expandafter\cmd@to@cs\string#1}}
\def\cmd@to@cs#1#2{\char\number`#2\relax}
\DeclareRobustCommand\cs[1]{\texttt{\char`\\#1}}

\newenvironment{Declaration}{%
%    \end{macrocode}
% \begin{macro}{\new@element}
%   Help macro to define new Declaration elements.
%    \begin{macrocode}
  \newcommand*{\new@element}[1]{%
    \expandafter\newcommand\expandafter*\csname X##1\endcsname{}%
    \expandafter\let\csname X##1\expandafter\endcsname
    \csname ##1\endcsname
    \expandafter\newcommand\expandafter*\csname new##1\endcsname[1]{%
%      \begingroup
%        \let\ensuremath\@firstofone
%        \let\textit\@firstofone
%        \lowercase{\def\@tempa{##1}}%
%        \pdfstringdef\@tempb{\label@base.\@tempa.####1}%
%        \xdef\@currentHref{\@tempb}%
%        \Hy@raisedlink{\hyper@anchorstart{\@currentHref}\hyper@anchorend}%
%        \label{desc:\label@base.\@tempa.####1}%
%      \endgroup
      \csname X##1\endcsname{####1}\ignorespaces
    }%
    \expandafter\let\csname ##1\expandafter\endcsname\csname new##1\endcsname
  }%
  \newcommand*{\new@xelement}[2]{%
    \expandafter\newcommand\expandafter*\csname X##1\endcsname{}%
    \expandafter\let\csname X##1\expandafter\endcsname
    \csname ##1\endcsname
    \expandafter\newcommand\expandafter*\csname new##1\endcsname[2]{%
%      \begingroup
%        \let\ensuremath\@firstofone
%        \let\textit\@firstofone
%        \lowercase{\def\@tempa{##1}}%
%        \pdfstringdef\@tempb{\label@base.\@tempa.####1.####2}%
%        \xdef\@currentHref{\@tempb}%
%        \Hy@raisedlink{\hyper@anchorstart{\@currentHref}\hyper@anchorend}%
%        \label{desc:\label@base.\@tempa.####1.####2}%
%      \endgroup
      \csname X##1\endcsname{####1}{##2{####2}}\ignorespaces
    }%
    \expandafter\let\csname ##1\expandafter\endcsname\csname new##1\endcsname
  }%
%    \end{macrocode}
%    \begin{macrocode}
  \new@element{Option}%
  \new@element{Macro}%
  \new@element{Environment}%
  \new@element{Counter}%
  \new@element{FloatStyle}%
  \new@element{PLength}%
  \new@element{Variable}%
  \new@xelement{OptionValue}{\PValue}%
%    \end{macrocode}
% \end{macro}
%    \begin{macrocode}
  \ifvmode\else\par\fi\addvspace{2\baselineskip}%
  \vspace{-\baselineskip}%
  \vspace{\z@ plus \baselineskip}%
  \noindent
  \start@Declaration
  \tabular{|l|}\hline\ignorespaces
}{%
  \\\hline\endtabular\nobreak\after@Declaration\nobreak\par\nobreak
  \vspace{1.5\baselineskip}\nobreak\vspace{-\baselineskip}\nobreak%
  \vspace{0pt minus .5\baselineskip}\nobreak%
  \aftergroup\@afterindentfalse\aftergroup\@afterheading
}
\newcommand*{\start@Declaration}{\hspace{-1em}}
\newcommand*{\after@Declaration}{}
% \begin{macro}{\Macro}
% \begin{macro}{\Option}
% \begin{macro}{\KOption}
% \begin{macro}{\OptionValue}
% \begin{macro}{\Environment}
% \begin{macro}{\Counter}
% \begin{macro}{\Length}
% \begin{macro}{\PLength}
% \begin{macro}{\FloatStyle}
% \begin{macro}{\Pagestyle}
% \begin{macro}{\Variable}
% \begin{macro}{\FontElement}
% \begin{macro}{\PName}
% \begin{macro}{\PValue}
% \begin{macro}{\Parameter}
% \begin{macro}{\OParameter}
% \begin{macro}{\AParameter}
% \begin{macro}{\PParameter}
% \begin{macro}{\POParameter}
%   \begin{description}
%   \item[\cs{Macro}] \LaTeX{} or \TeX{} macro
%   \item[\cs{Option}] class or package option
%   \item[\cs{KOption}] |\KOMAoptions| option
%   \item[\cs{Environment}] \LaTeX{} environment
%   \item[\cs{Counter}] \LaTeX{} counter
%   \item[\cs{Length}] \LaTeX{} length
%   \item[\cs{PLength}] \KOMAScript{} pseudo length
%   \item[\cs{Variable}] \KOMAScript{} variable
%   \item[\cs{FontElement}] \KOMAScript{} element that has its own font
%     selection
%   \item[\cs{PName}] name of a parameter of a macro or environment
%   \item[\cs{PValue}] value of a parameter of a macro or environment
%   \item[\cs{Parameter}] the mandatory parameter of a macro or environment
%   \item[\cs{OParameter}] the optional parameter of a macro or environment
%   \item[\cs{AParameter}] the alternativ parameter of a macro or environment
%   \item[\cs{PParameter}] the part-of-command parameter of a macro or
%     environment
%   \end{description}
%    \begin{macrocode}
\DeclareRobustCommand*{\Macro}[1]{\mbox{\texttt{\char`\\#1}}}
\DeclareRobustCommand*{\Option}[1]{\mbox{\texttt{#1}}}
\DeclareRobustCommand*{\KOption}[1]{\mbox{\Option{#1}\texttt=}}
\DeclareRobustCommand*{\OptionValue}[2]{\mbox{\texttt{#1=#2}}}
\DeclareRobustCommand*{\FloatStyle}[1]{\mbox{\texttt{#1}}}
\DeclareRobustCommand*{\Pagestyle}[1]{\mbox{\texttt{#1}}}
\DeclareRobustCommand*{\Environment}[1]{\mbox{\texttt{#1}}}
\DeclareRobustCommand*{\Counter}[1]{\mbox{\texttt{#1}}}
\DeclareRobustCommand*{\Length}[1]{\mbox{\texttt{\char`\\#1}}}
\DeclareRobustCommand*{\PLength}[1]{\mbox{\PValue{#1}}}
\DeclareRobustCommand*{\Variable}[1]{\mbox{\PValue{#1}}}
\DeclareRobustCommand*{\FontElement}[1]{\PValue{#1}}
\DeclareRobustCommand*{\PName}[1]{\texttt{\textit{#1}}}
\DeclareRobustCommand*{\PValue}[1]{\texttt{#1}}
\DeclareRobustCommand*{\Parameter}[1]{\texttt{\{}\PName{#1}\texttt{\}}}
\DeclareRobustCommand*{\OParameter}[1]{%
  \texttt{[%]
  }\PName{#1}\texttt{%[
    ]}}
\DeclareRobustCommand*{\AParameter}[1]{%
  \texttt{(%)
  }\PName{#1}\texttt{%(
    )}}
\DeclareRobustCommand*{\PParameter}[1]{\texttt{\{#1\}}}
\DeclareRobustCommand*{\POParameter}[1]{\texttt{[#1]}}
%    \end{macrocode}
% \end{macro}
% \end{macro}
% \end{macro}
% \end{macro}
% \end{macro}
% \end{macro}
% \end{macro}
% \end{macro}
% \end{macro}
% \end{macro}
% \end{macro}
% \end{macro}
% \end{macro}
% \end{macro}
% \end{macro}
% \end{macro}
% \end{macro}
% \end{macro}
% \end{macro}
% NOTE: taken from scrguide.cls

% \begin{environment}{desctable}
%   This is almost the same like \texttt{desctabular} but it uses a longtable
%   to allow page breaks.
%    \begin{macrocode}
\newenvironment{desctable}[1][2em]{%
  \onelinecaptionsfalse
  \start@desctab{#1}%
  \newcommand{\Endfirsthead}{\toprule\endfirsthead}%
  \newcommand{\Endhead}{\midrule\endhead}%
  \newcommand*{\standardfoot}{%
    \addlinespace[-.5\normalbaselineskip]\midrule
    \multicolumn{2}{r@{}}{\dots}\\
    \endfoot
    \addlinespace[-.5\normalbaselineskip]\bottomrule
    \endlastfoot
  }%
  \longtable{lp{\descwidth}}%
}{%
  \endlongtable
}
%    \end{macrocode}
% \end{environment}

% \begin{length}{\descwidth}
%   I need a length of local usage. I could have used |\@tempdima| or
%   another local length from kernel. But I've decided not to try to find a
%   unused length at \texttt{tabular} environment.
%    \begin{macrocode}
\newlength{\descwidth}
%    \end{macrocode}
% \end{length}

% \begin{macro}{\start@desctab}
%   This is the \emph{worker} macro of \texttt{desctable} and
%   \texttt{desctabular}. It does the complete calculations and definition of
%   the entry (something like |\item|) commands.
%    \begin{macrocode}
\newcommand*{\start@desctab}[1]{%
  \setlength{\descwidth}{\linewidth}%
  \addtolength{\descwidth}{-4\tabcolsep}%
  \addtolength{\descwidth}{-#1}%
  \setlength{\labelwidth}{\linewidth}%
  \addtolength{\labelwidth}{-2\tabcolsep}%
  \newcommand{\nentry}[2]{%
    \multicolumn{2}{p{\labelwidth}}{\raggedright##1}\\*%
    \hspace*{#1} & ##2\tabularnewline%
  }%
  \newcommand{\entry}[2]{\nentry{##1}{##2}[.5\baselineskip]}%
  \newcommand*{\pentry}[1]{%
    \entry{\PLength{##1}\IndexPLength[indexmain]{##1}}}%
  \newcommand*{\pventry}[1]{\entry{\PValue{##1}}}%
  \newcommand*{\mentry}[1]{\entry{\Macro{##1}}}%
  \newcommand*{\ventry}[1]{%
    \entry{\Variable{##1}%\IndexVariable[indexmain]{##1}%
    }%
  }%
  \newcommand*{\feentry}[1]{%
    \entry{\FontElement{##1}\IndexFontElement[indexmain]{##1}}%
  }%
  \newcommand*{\oentry}[1]{%
    \entry{\Option{##1}\IndexOption[indexmain]{##1}}%
  }%
}
% \end{macro}

\makeatother

\newcommand{\tubslatex}{\emph{tubslatex}\xspace}

\def\example{\par\smallskip\noindent\textit{Beispiel: }}

% Schreibt 'Beispiel vor den folgenden Inhalt und rückt alles nach dem ersten
% Absatz um 2em ein.
\newenvironment{Example}{%
\begingroup
\leftskip2em
\par\smallskip\noindent\hspace*{-2em}\textit{Beispiel: }
}{%
\par\endgroup
}

% Umgebung 'Wichtig:' ...
\newenvironment{important}{%
  \begin{description}
    \item[\itshape\mdseries\rmfamily Wichtig:]
}{%
  \end{description}
}

\newcommand{\zB}{\mbox{z.\,B.}\xspace}


%opening
\title{\tubslatex -- Das CorporateDesign in \LaTeX}
\subtitle{Anleitung und Dokumentation}
\author{Enrico Jörns}
\publishers{Presse und Kommunikation} % to logo?

\begin{document}


%Befehle für Glossar
\newglossaryentry{glos:siegelbandlogo}{%
  name=Siegelbandlogo,
  description={\tubslogo}
}
\newglossaryentry{glos:gaussraster}{%
  name={Gau\ss raster},
  description={Auf der \glslink{glos:summenformel}{gauß'schen Summenformel} basierende Unterteilung der Seite
    in Segmente. Benachbarte Segmente können beliebig zusammen gefasst werden}
}
\newglossaryentry{glos:modulsystem}{%
  name={Modulsystem},
  description={Flexibles Platzierungssystem für wissenschaftliche Plakate.
  Dabei wird der Darstellungsbereich komplett in einzelne Module
  verschiedener Größe aufgeteilt.}
}
\newglossaryentry{glos:cmyk}{%
  name={CMYK-Farbmodell},
  description={Das CMYK-Farbmodell ist ein subtraktives Farbmodell,
    das die technische Grundlage für den modernen Vierfarbdruck bildet.
    Die Abkürzung CMYK steht für die drei Farbbestandteile
    \emph{Cyan}, \emph{Magenta}, \emph{Yellow}
    und den Schwarzanteil \emph{Key} als Farbtiefe}
}
\newglossaryentry{glos:absenderbereich}{%
  name={Absenderbereich},
  description={
    Freier Bereich am oberen oder unteren Blattrand
    zur Darstellung eines Absenders (Institut/zentrale Eintrichtung).
    Ihm schließt sich direkt der \gls{glos:kommunikationsbereich} mit dem Inhalt an.\\
    Die Position des Absenderbereichs kontrolliert beim Gauß-Layout auch die
    Reihenfolge der Segmentaufteilung. Das größte Segment befindet sich immer
    auf der Seite des Absenderbereichs.
    }
}
\newglossaryentry{glos:kommunikationsbereich}{%
  name={Kommunikationsbereich},
  description={Bereich zur Darstellung von Inhalten}
}
\newglossaryentry{glos:spaltenraster}{%
  name={Spaltenraster},
  description={Abhängig vom Format kann der Textbereich einer Seite in
    6, 4 oder 2 Grundspalten geteilt werden, welche alle die selbe Breite und
    den selben Abstand zueinander haben.
    Benachbarte Grundspalten können variabel zu einer Darstellungsspalte
    zusammengefasst werden}
}
\newglossaryentry{glos:bindekorrektur}{%
  name={Bindekorrektur},
  description={}
}
\newglossaryentry{glos:sekundaerfarbklang}{%
  name={Sekund\"arfarbklang},
  description={}
}
\newglossaryentry{glos:mediaevalziffern}{%
  name={Medi\"avalziffern},
  description={Ziffern, die im Gegensatz zu Versalziffern Ober- und Unterlänge
  haben und sich dadurch im Mengentext besser in das Schriftbild einfügen als \gls{glos:versalziffern}}
}
\newglossaryentry{glos:versalziffern}{%
  name={Versalziffern},
  description={Auf der Grundlinie ausgerichtete Ziffern ohne Ober- und Unterlängen.
  Sie fügen sich daher meist schlechter in Schriftbild ein als \gls{glos:mediaevalziffern}}
}
\newglossaryentry{glos:dinlang}{%
  name={DIN lang},
  description={Bezeichnet hier ein Format von $1/3$ der Blattgröße \mbox{DIN\,A4},
  wie es zum Beispiel für Flyer (doppelt gefaltet) häufig verwendet wird.}
}
\newglossaryentry{glos:summenformel}{%
  name={Gauß'sche Summenformel},
  description={Berechnungsgrundlage für die formatübergreifende
  vertikale Aufteilung des Kommunikationsbereichs im CD (siehe \gls{glos:gaussraster}).
  \[1 + 2 + 3 + 4 + \ldots + n = \sum_{k=1}^n k = \frac{n(n+1)}{2}\]}
}

%Akronyme
\newacronym{CD}{CD}{Corporate Design}
\newglossaryentry{CMYK}{%
  type=\acronymtype,
  name={CMYK},
  description={Cyan, Magenta, Yellow, Key (siehe \gls{glos:cmyk})}
}
\newacronym{RGB}{RBG}{Red, Green, Blue}



\pagestyle{empty}
\hfill
\begin{center}
\huge\color{tuRed}
Dokument befindet sich im Aufbau!\bigskip

Informationen evtl. unvollständig oder unklar!
\end{center}
\vspace*{3cm}
\begin{center}
\LARGE\tubslatexVersion
\end{center}

\newpage~\newpage

\begin{titlepage}
  \showtubslogo
  \showlogo{Presse und Kommunikation}
  \showtopline
  \begin{segment}{2}
    \centering
    \Huge \tubslatex\\[\medskipamount]
    \bfseries\huge Das Corporate Design in \LaTeX
  \end{segment}
  \begin{segment}{1}
    \centering
    \LARGE Anleitung und Dokumentation
  \end{segment}
  \begin{segment}{2}
    \centering
    \Large Enrico Jörns\\[\bigskipamount]
    \bfseries\today
  \end{segment}
  \begin{segment}{3}
    \Large \tubslatex-Autoren: Enrico Jörns, Tobias Rad, Martin Bäker et al.
  \end{segment}
\end{titlepage}

\begin{backpage}
  \showtubslogo
  \showtopline
  \begin{segment}{8}
    \begin{tabular}{@{}ll@{}}
      Version:          & \tubslatexVersion \\
      Projektseite:     & http://enricojoerns.de/tubslatex \\
      Projektleitung:   & TU Braunschweig, Presse und Kommunikation \\
      Kontakt:          & Enrico Jörns -- e.joerns@tu-bs.de \\
    \end{tabular}
    \vfill
    (c) 2012 by TU Braunschweig
  \end{segment}
\end{backpage}


\chapter*{Vorwort}

Eine klare, wiedererkennbare und dennoch abwechslungsreiche
Repräsentation ist in einem medialen Umfeld wie wir es heutzutage
vorfinden nicht nur für jedes Unternehmen, sondern vor allem auch
für eine Universität von großer Bedeutung.
Sie schafft Identität, fördert Identifikation und behauptet sich zugleich im 
Blickfeld der Öffentlichkeit

Mit dem Corporate Design der TU Braunschweig ist dieses wichtige Ziel
erfüllt worden. Es verbindet ein ganzheitliches Aussehen mit dem nötigen
Spielraum für Individualität und Kreativtät der einzelnen Institute und 
Einrichtungen.\bigskip


Die Vorgaben und Möglichkeiten, die das Corporate Design bietet,
mit einem effektiven und professionellen Textsatz zu verbinden, ist
das Ziel von \tubslatex. Diese Sammlung von Paketen und Klassen
für das Textsatzprogramm \LaTeX\ soll dem Benutzer zum einen das Erstellen und
die Gestaltung von verschiedensten Schriftstücken vereinfachen und zum anderen
die Einhaltung der notwendigen Richtlinien sicherstellen.

Die Einsatzmöglichkeiten sind dabei sehr vielfältig.
Vom Schrieben einfacher Textdokumente und Briefe über die Gestaltung von Postern
und Flyern bis hin zur Erstellung von umfangreichen Präsentationen erstreckt
sich das Spektrum der Möglichkeiten.
Im Mittelpunkt fast jedes Dokumentes stehen die Grundelemente des Corporate
Design. Dazu zählen das rote Siegelbandlogo, die Haussschrift Nexus,
das markante Gauß-Layout, sowie ein klares Farbschema.\bigskip


Alle Vorlagen wurden nach bestem Wissen erstellt und getestet, was
jedoch den Fehlerteufel erfahrungsgemäß zumeist nicht sonderlich beeindruckt.
Daher bitte ich bei Auftreten eines Fehlers oder Problemes Milde walten
zu lassen, den beherzten Schritt zu einem Mail-Programm Ihres Vertrauens zu
wagen und den Mangel kund zu tun. Dies gilt natürlich auch für Fragen,
Anregungen und Wünsche.

Die aktuelle E-Mail-Adresse dafür lautet: \textit{e.joerns@tu-bs.de}.

\hfill\textit{Braunschweig, den \today}\\[\bigskipamount]
\noindent Enrico Jörns



\pagestyle{scrheadings}% TODO:!?!?

\phantomsection % requires hyperref package
\pdfbookmark[1]{Inhaltsverzeichnis}{toc}
\tableofcontents


\newcommand{\newdocumentclass}[1]{\textcolor{tuRed}{\lstinline{#1}}}
\newcommand{\newpackage}[1]{\textcolor{tuRed}{\lstinline{#1}}}

\chapter{Einleitung}

Das vorliegende Dokument beschreibt die Verwendung die Möglichkeiten von \tubslatex.
Mit \tubslatex\ ist im Folgenden immer die Sammlung aller Vorlagen in Form
von \LaTeX-Paketen und -Klassen bezeichnet, welche für die Darstellung von
Dokumenten im Corporate Design der TU Braunschweig vorgesehen sind.

Das folgende Kapitel stellt zuerst die wesentlichen Elemente und Richtlinien
des Corporate Designs vor und ... % TODO...

Das Kapitel~\ref{chap:install} bietet einige kleine Hinweise, die bei
der Installation von \tubslatex nützlich sein könnten.
Es beschreibt außerdem Möglichkeiten zur manuellen Installation, falls
kein Installer für das gewünschte Betriebssystem installiert oder eine
individuelle Installation durchgeführt werden soll.

Der erste Teil des Dokumentes beschäftigt sich mit der allgemeinen Verwendung
der vorgestellten Klassen.
Dabei wurde versucht, für jede Klasse eine möglichst vollständige Beschreibung
zu liefern. Da sich allerdings viele Funktionen gleichen, das sie auf die selben
Grundfunktionen zurück greifen sind manche Funktionen an einer Stelle nur kurz
erwähnt mit einem Verweis auf die ausführlichere Beschreibung.

Einige Grundelemente von \tubslatex sind in eigenständigen Paketen untergebracht,
sodass sie bei Bedarf auch von anderen \LaTeX-Klassen benutzt werden könnten.
Dazu gehören unter anderem die Schriftart \emph{Nexus}, das Siegelband-Logo
und die Farbdefinitionen.
Auf diesen Aspekt geht der zweite Teil des Dokumentes näher ein.% TODO: Link

Für besonders Eilige empfiehlt sich das Kapitel~\ref{chap:rapid},
welches in aller Kürze die jeweils wichtigsten Befehle und Optionen einer Klasse
zusammenfasst und so einen raschen Einstieg ermöglicht.

\clearpage
\section{Das Corporate Design der TU Braunschweig}

\paragraph{Seitenlayout}\hfill\\
\begin{minipage}[t]{0.45\textwidth}
\vspace*{0pt}
\centering\fboxsep0mm%
\fbox{\includegraphics[width=0.85\textwidth,page=1]{examples/thecd.pdf}}
\end{minipage}
\begin{minipage}[t]{0.55\textwidth}
\vspace*{0pt}
Das Corporate Design setzt sich aus ein paar wesentlichen Elementen und
Gestaltungsmerkmalen zusammen.
Eine Seite teilt sich grundsätzlich in einen \emph{Absenderbereich} (1) und einen
\emph{Kommunikationsbereich} (2).
Das markante rote \emph{Siegelband-Logo} dient dabei als Bindeglied zwischen
diesen beiden Bereichen. 
Der Absenderbereich kann sowohl am oberen als auch am unteren Blattende
platziert werden.
Das Siegelband-Logo passt sich in seiner vertikalen Position an und kann darüber hinaus
jeweils rechts- oder linksseitig platziert werden.
Daraus ergeben sich insgesamt 4 Darstellungsvarianten.

Der Kommunikationsbereich wird von einem Rand fester Breite eingerahmt,
der der Gestaltung einen edlen Wiedererkennungswert verleiht und gleichzeitig
den Inhalt absetzt, um Druck- und Kopierprobleme zu vermeiden.
\end{minipage}\bigskip

\paragraph{Absenderfeld und Gaußraster}\hfill\\
\begin{minipage}[t]{0.45\textwidth}
\vspace*{0pt}
\centering\fboxsep0mm%
\fbox{\includegraphics[width=0.85\textwidth,page=2]{examples/thecd.pdf}}
\end{minipage}
\begin{minipage}[t]{0.55\textwidth}
\vspace*{0pt}
Das im Absenderbereich befindliche \emph{Absenderfeld} (1) bietet Platz für
einen Instituts-/Abteilungsnamen bzw. ein entsprechendes Logo.

Der Kommunikationsbereich kann in einem an die gaußsche Summenformel
angelehnten Raster (\emph{Gaußraster}) unterteilt werden.
Dabei ergibt sich die Hohe eines Grundsegmentes ((2), (3), \ldots)
immer aus der Summer der Höhe der beiden vorhergehenden Grundsegmente.
Zur Darstellung können beliebig viele Grundsegmente zusammengefasst werden,
womit sich ein zugleich flexibles aber trotzdem charakteristisches
Gesamtbild ergibt.

Das Gaußraster ist immer so ausgerichtet, dass sich das breiteste Grundsegment
an der Grenze zum Absenderbereich befindet.
\end{minipage}

\paragraph{Schrift und Farbwelt}\hfill\\
\begin{minipage}[t]{0.45\textwidth}
\vspace*{0pt}
\centering\fboxsep0mm%
\fbox{\includegraphics[width=0.85\textwidth,page=3]{examples/thecd.pdf}}
\end{minipage}
\begin{minipage}[t]{0.55\textwidth}
\vspace*{0pt}
Zum Umfang des Corporate Design gehört neben den verschiedenen
Gestaltungsmerkmalen auch eine unverkennbare Schrift.
Dieses ist mit der modernen Hausschrift \emph{Nexus} gegeben, die sowohl mit als auch
ohne Serifen in jeweils drei Schriftschnitten zur Verfügung steht.

Ein weiteres wichtiges Merkmal von Dokumenten im Corporate Design sind
die verwendeten Farben, die sich aus einer Reihe fest definierter Paletten
zusammen setzen.
Als wichtige Strukturfarben sind das TU-Rot, sowie Schwarz und Weiß vorgesehen.
Zur individuellen Gestaltung stehen 4 Farbklänge (Gelb/Orange, Grün, Blau, Violett)
mit jeweils 3 Farben und 12 Abstufungen zur Verfügung.
\end{minipage}

\clearpage
\section{Übersicht über die Vorlagen}

Die Vorlagen bauen allesamt auf Standard-\LaTeX-Klassen und Paketen auf.
Insbesondere sind alle Dokumentenklassen von Klassen des KOMA-Skripts
abgeleitet, was sich auch in ihrem Benennungsschema widerspiegelt.
Konkret bauen die Vorlagen für Poster (\texttt{tubsposter}) und Dokumente
(\texttt{tubsartcl}, \texttt{tubsreprt}, \texttt{tubsboox}) auf
den entsprechenden KOMA-Skript-Klassen auf,%
\footnote{\texttt{tubsposter} verwendet \texttt{scrartcl}}.
Dabei wurden einige Optionen und Befehle, die die Basis-Klassen bieten, erhalten
und können in der ausführlichen Dokumentation zum KOMA-Skript\cite{koma-skript}
nachgeschlagen werden.1
Einige Funktionalitäten wie etwa die Satzspiegelberechnung sind dagegen
komplett ersetzt wurden.
Die Dokumentation geht an einigen Stellen noch expliziter auf Unterschiede ein.

Ebenfalls von einer KOMA-Klasse abgeleitet ist die Briefklasse
\texttt{tubslttr2}.
Diese nutzt (wie der Name vermuten lässt) die noch verhältnismäßig junge
Klasse \texttt{scrlttr2}. Dies spiegelt sich vor allem darin wieder, dass sich
das verwendete Interface in Teilen stark von den anderen Klassen unterscheidet.
Die Dokumentation geht dabei hauptsächlich auf die Besonderheiten in
\texttt{tubslatex} ein.
Für detailliertere Informationen kann hier auch wieder die Dokumentation von
KOMA\cite{koma-skript} zu Rate gezogen werden.

Für Präsentationen wird die relativ bekannte \texttt{beamer}-Klasse verwendet.
Hier stellt \tubslatex keine eigene Klasse zur Verfügung,
sondern bietet entsprechend der Beamer-Philosophie ein \emph{Beamer-Theme} an,
das im Dokument geladen werden kann und so viel Flexibilität erlaubt.
Bei der Verwendung ergeben sich daher lediglich eine minimale Anzahl an
Veränderungen und Erweiterungen des Funktionsumfangs von beamer.

Die Vorlage für Broschüren (\texttt{tubsleaflet}) ist
von der Klasse \texttt{leaflet} abgeleitet.


       % Kapitel 'Einleitung'

\chapter{Installation}\label{chap:install}

Aktuell wird \tubslatex in drei verschiedenen Varianten angeboten.
Zum einen als gepackter tex-Paketaum, der sowohl lokal als auch systemweit
entpackt werden kann. Allerdings müssen hier alle notwendigen Schritte zur 
Installation der Schriften und Anmelden der neuen Datein manuell gemacht werden.
Für Windows-Benutzer, die MikTeX verwenden, gibt es einen automatischen
Installer, der alle nötigen Schritte automatisch ausführt.
Ebenso gibt es für Debian-basierte Systeme, die texlive verwenden ein
\texttt{.deb}-Paket, das alle Abhängigkeiten korrekt auflöst und die Vorlagen
automatisch installiert.

\section{Automatische Installation}

\subsection{Windows -- MikTeX}

Bei Verwendung von MikTeX unter Windows kann auf den verfügbaren Installer
zurück gegriffen werden. Um eine fehlerfreie Installation zu ermöglichen
sollte MiKTeX in der Version 2.9 vorhanden sein.
Ältere Versionen werden nur bedingt unterstützt.

Der Installer kopiert alle benötigten Dateien in ein frei wählbares Verzeichnis
auf der Festplatte und registriert dies automatisch als neuen texmf-Baum.
Auch die Schriftart Nexus wird standardmäßig installiert und in 
Font-Verwaltung registriert.

\subsection{Ubuntu/Debian -- TexLive}

Für debian-basierte Systeme wie Ubuntu oder Debian selber steht ein
\texttt{.deb}-Paket zur Verfügung, dass einfach mit Hilfe eines Paketmanagers
installiert werden kann.
Genutz werden können dafür die diversen graphischen Tools oder einfach das
Kommandozeilenprogramm \texttt{gdebi}. Diese Programme lösen alle Abhängigkeiten
auf und installieren sie automatisch mit.

Die Installation der Dateien erfolgt in das Verzeichnis
\lstinline{/usr/share/texmf-texlive/}. Sowohl die neuen Dateien als auch 
die Schrift Nexus werden automatisch registriert.

\section{Manuelle Installation}

Zur manuellen Installation liegen die LaTeX-Dateien als gepacktes
zip-Verzeichnis vor, in dem die komplette tex-Pfadstruktur oberhalb der
texmf-Ebene enthalten ist. Konkret sind dies die Ordner
\lstinline{doc}, \lstinline{tex} und \lstinline{fonts}.

Diese können entweder lokal oder global entpackt werden.
Für erste Variante kann der Inahlt einfach in das Verzeichnis
\lstinline{/home/benutzername/texmf} hinein geschrieben werden.

Es sei darauf hingewiesen, dass alle weiteren benötigten Pakete der
Tex-Distribution ggf. manuell installiert werden müssen.
MikTeX macht dies normalerweise automatisch.
Bei texlive können unter Linux zum Beispiel die entsprechenden System-Pakete
installiert werden.


\subsection{MikTeX}

\begin{description}
  \item[1. Dateien Kopieren]
  \item[2. Dateien registrieren]
  \item[3. Schriften registrieren]
\end{description}


\subsection{TexLive}

\begin{description}
  \item[1. Dateien Kopieren] Zuerst müssen alle in der Zip-Datei enthaltenen
    Dateien entpackt werden.
    Diese liegen dort schon in der korrekten Ordner-Sktruktur vor, sodass
    sie lediglich in das Ziel-texmf-Verzeichnis \emph{hinein}kopiert werden
    müssen.
    
    Grundsätzlich kann man sich bei der Installation entscheiden,
    ob die Vorlagen nur Lokal für den aktuellen Benutzer oder systemweit
    installiert werden sollen.
    Es sei darauf hingewiesen, dass eine lokale Installation ein paar
    Fallstricke mit sich bringt, die evtl. zu Problemen bei der späteren
    Installation weiterer Pakete führen können.
    
    Entpacken als Benutzer ins lokale texmf-Verzeichnis:
    \begin{lstlisting}
$ unzip -d ~/texmf/ tubslatex_0.3-alpha2.zip
    \end{lstlisting}
    Oder als Root ins systemweite texmf-Verzeichnis:
    \begin{lstlisting}
# unzip -d /usr/share/texmf/ tubslatex_0.3-alpha2.zip
    \end{lstlisting}

  \item[2. Dateien registrieren]
  \item[3. Schriften registrieren]\hfill

    \begin{lstlisting}
# /usr/bin/updmap-sys --nomkmap --nohash --enable Map=NexusProSans.map
# /usr/bin/updmap-sys --nomkmap --nohash --enable Map=NexusProSerif.map
# /usr/bin/updmap-sys
    \end{lstlisting}
    
    Für eine lokale Installation muss jeweils der Zustatz \texttt{-sys}
    weggelassen werden.

\end{description}


\newenvironment{knownissue}[1]{%
  \paragraph{#1}\hfill
  \newcommand{\solution}[1]{\noindent{\itshape ##1}}
}{%
}

\section{Fehlerbehebung}

\subsection{MikTeX}

\subsection{texlive}

\begin{knownissue}{%
  updmap bricht plötzlich ab (Fehlercode 2) / getnonfreefonts bricht ab
}

Dieses Problem ist unter anderem von texlive2009 bekannt.
Dort hat sich ein Problem eingeschlichen, das leider die Funktionalität von
updmap teilweise lahmlegt.

\solution{Mögliche Lösung:}
\begin{itemize}
  \item Die Datei 10local.cfg (falls sie noch nicht vorhanden ist) manuell
    erstellen (z.B. mit \texttt{touch}).
  \item Nicht die Option \texttt{--quiet} verwenden.
\end{itemize}

Alternativ kann auch die Zeile \lstinline{set -e} am Anfang des Skripts einfach auskommentiert werden.
\end{knownissue}


\begin{knownissue}{
  \texttt{mktexpk: don't know how to create bitmap font for \ldots}\newline
  nach Schriften-Installation mit \lstinline{updmap-sys}}

Dieses Problem kann auftauchen, wenn auf dem System vom aktuellen Benutzer
zuvor schon einmal \lstinline{updmap} (ohne -sys) ausgeführt wurde.
Dann legt texlive eine lokale Map-Datei an und verwendet fortan diese,
auch wenn die System-Map mit \lstinline{updmap-sys} upgedated wurde.

\noindent\textit{Überprüfung:}

Die Ausgabe von
\begin{lstlisting}
kpsewhich pdftex.map
\end{lstlisting}
sollte einen anderen Pfad als
\lstinline{/var/lib/texmf/fonts/map/pdftex/updmap/} liefern.

\solution{Lösung 1:}

Entfernen der aktuell verwendeten (lokalen) Map-Datei.

\solution{Lösung 2:}

Installation weiterer Schriften mit \lstinline{updmap} statt
\lstinline{updmap-sys}.

\end{knownissue}




% TODO...

\def\foonote{%
kpathsea: Running mktexpk --mfmode / --bdpi 600 --mag 1+240/600 --dpi 840 NexusProSans-Bold-Regular--base
mktexpk: don't know how to create bitmap font for NexusProSans-Bold-Regular--base.

Reproduce:
  - sudo updmap
  - sudo aptitude install texlive-tubs
Test:
  - wenn 'kpsewhich pdftex.map' nicht /var/lib/.. ausgibt, liegt der Fehler vor
Beschreibung:
  - Dieses Problem kann unter anderem auftreten, wenn vor der Installation
    irgendwann einmal 'updmap' aufgerufen wurde.
Begründung:
  - Installation greift auf updmap-sys zurück, die Damit generierte Datei
    wird aber nicht mehr verwendet.
Lösung 1:
  - updmap manuell aufrufen.
Lösung 2:
  - Datei, die bei 'kpsewhich pdftex.map' angezeigt wird, löschen
    (z.B. /root/.texmf-var/fonts/map/pdftex/updmap/pdftex.map)
  - 'sudo updmap-sys' aufrufen
}
     % Kapitel 'Installation'

\chapter{Schnellstart}\label{chap:rapid}

% TODO: Schnelleinführung Corporate Desing (Gaußraster, etc.)

\paragraph{Dokumente}
Die schnellste Methode ein Dokument im Corporate Design zu erstellen ist
das Laden einer der zur Verfügung stehenden Dokumentenklassen.
Für die Erstellung von Textdokumenten sind dies \newdocumentclass{tubsartcl},
\newdocumentclass{tubsreprt} und \newdocumentclass{tubsbook}.

\paragraph{Poster}
Die Klasse \newdocumentclass{tubsposter} kann zum Erstellen von Postern verwendet werden.
% tubsflyer?

\paragraph{Papierformat}
Das zu verwendende Papierformat sollte dabei als optionales Argument mit
übergeben werden. Für ein Dokument in DIN A4 ist dies \texttt{a4paper}.
Zur Verfügung stehen alle Papierformate von A0 bis A6.

Standardmäßig wird die Titelseite in Dokumenten einfach mit dem Logo und einer
roten Trennlinie zwischen Kommunikations- und Absenderbereich versehen.

\paragraph{Briefe}

Die Klasse \newdocumentclass{tubslttr2} ist für die Erstellung von Briefen
vorgesehen.

\paragraph{Präsentationen}
Um Präsentation zu erstellen existiert ein Style für \LaTeX-Beamer.
Dieser wird einfach mit \lstinline!\usetheme{tubs}! geladen.
  % Kapitel 'Schnellstart'


\part{Anwenderklassen und -Pakete}

% \chapter{Dokumente}

Einfache Textdokumente können mit den Klassen \newdocumentclass{tubsartcl},
\newdocumentclass{tubsreprt} und \newdocumentclass{tubsbook} erstellt werden.

% \section{Bindekorrektur und Marginalen}

Da Dokumente normalerweise gedruckt werden und ggf. auch gebunden, werden
sie standardmäßig mit einer kleinen Bindekorrektur gesetzt, sodass das
Logo beim Drucken oder Abheften nicht abgeschnitten wird.
Diese Bindekorrektur kann mittels der Paketoption \texttt{bcor} angepasst werden.
Mit \texttt{bcor=0mm} wird sie beispielsweise deaktiviert.



\section{Titelseite}

Titelseiten können bei \LaTeX\ generell auf zwei verschiedene Arten erstellt
werden; entweder mit Hilfe des Befehls \lstinline{\maketitle} oder mit
der Umgebung \lstinline{titlepage}. Beide Varianten werden von \tubslatex
unterstützt und leicht modifiziert.

\begin{Declaration}
  \Macro{maketitle}\OParameter{style}
\end{Declaration}

Die einfache Verwendung von \lstinline{\maketitle} erzeugt eine Titelseite
mit dem TU-Logo und einer roten Trennlinie zwischen Absender- und
Kommunikationsbereich.

Mit Hilfe des optionalen Arguments \PName{style} kann die Darstellung
der Titelseite geändert werden, indem aus einer Reihe vordefinierter Styles
ausgewählt wird. Eine weiterführende Erklärung findet
sich dazu in Kapitel~\ref{sec:titlestyles}.\bigskip

Neben den standardmäßig definierten Elementen für Titelseiten wie \Macro{author}
oder \Macro{title} werden in \tubslatex noch ein paar Zusätzliche definiert.

\begin{Declaration}
  \Macro{logo}\Parameter{logo}\\
  \Macro{titlepicture}\Parameter{file}\\
  \Macro{titleabstract}\Parameter{text}
\end{Declaration}

\Macro{logo} dient zur Darstellung eines zusätzlichen Absenders als Schrift
oder Bild. Es wird in allen Stilen im Absenderbereich auf der dem
TU Siegelband-Logo gegenüberliegenden Seite dargestellt. Eingefügte Grafiken
werden automatisch auf die Höhe des Kopfbereichs skaliert.

Mit \Macro{titlepicture} kann eine Bilddatei angegeben werden, die bei
Verwendung eines entsprechenden Styles auf der Titelseite dargestellt wird.

Der Befehl \Macro{titeabstract} erlaubt die Darstellung eines kurzen zusammenfassenden Textes auf der Titelseite, sofern der gewählte Stil dies unterstützt.
Weitere Erläuterungen finden sich jeweils auch im Kapitel~\ref{sec:titlestyles}.


\begin{Declaration}
  \XMacro{begin}\PParameter{\Environment{titlepage}}\\
  \quad\dots\\
  \XMacro{end}\PParameter{titlepage}
\end{Declaration}

Mit der \Environment{titlepage}-Umgebung können, wie von den Standardklassen
gewohnt, Titelseiten definiert werden. Sie bietet aber in \tubslatex noch einige
zusätzliche Darstellungsmöglichkeiten.

\begin{Declaration}
  \Macro{showtubslogo}\OParameter{Seite}\\
  \Macro{showlogo}\Parameter{Inhalt}\\
  \Macro{showtopline}
\end{Declaration}

Die Verwendung einer dieser Befehle innerhalb einer Titlepage-Umgebung sorgt
dafür, dass das entsprechende Element auf der Seite dargestellt wird.
\Macro{showtubslogo} stellt das Siegelbandlogo dar. Der optionale Parameter
\PName{Seite} akzeptiert legt dabei fest auf welcher Seite des Blattes
das Logo gezeigt werden soll. Der Wert \PValue{left} entspricht der Standardeinstellung und stellt das Logo auf der linken Seite,
% bzw. bei zweiseitigem Layout auf der Innenseite % NOTE: für Titelseite irrelevant!?
dar.
Mit dem Wert \PValue{right} wird dagegen eine Darstellung auf der rechten
% bzw. äußeren
Seite des Blattes erwirkt. Gleichzeitig hat dies Auswirkungen
auf die Positionierung des Zweitlogos, das automatisch immer auf der 
dem Siegelbandlogo gegenüberliegenden Seite platziert wird.
% TODO: inner, outer?

\begin{Declaration}
  \XMacro{begin}\PParameter{\Environment{titlerow}}%
    \OParameter{options}%
    \Parameter{gaussheight}\\
  \quad\dots\\
  \XMacro{end}\PParameter{titlerow}
\end{Declaration}

Die Umgebung \Environment{titlerow} erlaubt es dabei, die Titelelemente im 
Gaußraster anzulegen. Der Parameter \PName{gaussheight} gibt dabei
die Höhe des jeweiligen Elements in Segmenten an. Die Position der Elemente
ergibt sich aus der Reihenfolge der Definition.
Mit dem optionalen Parameter \PName{options} können Einstellung wie die Hintergrundfarbe oder ein Hintergrundbild übergeben werden.

\subsection{Rückseiten}

Neben der Möglichkeit Titelseiten zu definieren, gibt es in \tubslatex auch
einen ähnlichen Mechanismus zur Erstellung von Rückseiten derselbigen, sodass
ohne Mühe ein komplettes Titelblatt, dass auf die Möglichkeiten des
Corporate Design abgestimmt ist, erstellt werden kann.

\begin{Declaration}
  \Macro{makebackpage}\OParameter{Style}
\end{Declaration}

Dieser Befehl stellst das Gegenstück zu \Macro{maketitle} auf und erstellt eine
Rückseite. Genau wie bei \Macro{maketitle} kann über den Optionalen Parameter
\PName{style} wieder ein vordefinierter Stil geladen werden.

Außerdem stehen ein paare Befehle zur Verfügung, mit denen sich der Inhalt
der Rückseite festlegen lässt. Die konkrete Darstellung ist dabei wieder
von dem verwendeten Style abhängig. Weiteres findet sich in 
Kapitel~\ref{sec:backstyles}.


\begin{Declaration}
  \Macro{address}\Parameter{Adressdaten}\\
  \Macro{backpageinfo}\Parameter{Inhalt}
\end{Declaration}

Mit \Macro{address} kann eine Herausgeber-Adresse für das Dokument
dargestellt werden.
mit \Macro{backpageinfo} Lässt sich ganz allgemein der verbleibende
Inhalt der Rückseite nach Belieben füllen.
Siehe hierzu jeweils auch Kapitel~\ref{sec:backstyles}.

\subsection{Vordefinierte Titel-Styles}\label{sec:titlestyles}

Es sind 3 einfache Styles definiert.% vordefiniert
% TODO: styles selbst definieren.
% Hinweis zu Farbdarstellung!
% TODO: Default-Style in Config-Datei speichern

\begin{center}
  \fboxsep0mm
  \begin{minipage}[t]{0.33\textwidth}
    \centering\sffamily
    \fbox{%
      \includegraphics[width=0.95\textwidth]{examples/titlestyle_plain.pdf}}
    [default]
  \end{minipage}%
  \begin{minipage}[t]{0.33\textwidth}
    \centering\sffamily
    \fbox{%
      \includegraphics[width=0.95\textwidth]{examples/titlestyle_image.pdf}}
    [image]
  \end{minipage}%
  \begin{minipage}[t]{0.33\textwidth}
    \centering\sffamily
    \fbox{%
      \includegraphics[width=0.95\textwidth]{examples/titlestyle_imagetext.pdf}}
    [imagetext]
  \end{minipage}
\end{center}

\subsection{Vordefinierte Rückseiten-Styles}\label{sec:backstyles}

Es sind 3 einfache Styles definiert.% vordefiniert

\begin{center}
  \fboxsep0mm
  \begin{minipage}[t]{0.33\textwidth}
    \centering\sffamily
    \fbox{%
      \includegraphics[width=0.95\textwidth]{examples/backstyle_plain.pdf}}
    [plain] (default)
  \end{minipage}%
  \begin{minipage}[t]{0.33\textwidth}
    \centering\sffamily
    \fbox{%
      \includegraphics[width=0.95\textwidth]{examples/backstyle_info.pdf}}
    [info]
  \end{minipage}%
  \begin{minipage}[t]{0.33\textwidth}
    \centering\sffamily
    \fbox{%
      \includegraphics[width=0.95\textwidth]{examples/backstyle_addressinfo.pdf}}
    [addressinfo]
  \end{minipage}
\end{center}

\section{Kopf-/ Fußzeile}

Standardmäßig wird die Fußzeile komplett leer gelassen und die Kopfzeile
wird mit Seitennummer rechts und Kapitelname links gesetzt. Bei zweiseitigem
Layout gilt dies für die ungeraden Seiten, gerade Seiten werden entsprechend
mit Seitennummer links und Kapitelname rechts gesetzt.

\begin{Declaration}
  \Macro{ihead}\Parameter{innen}\\
  \Macro{ohead}\Parameter{außen}
\end{Declaration}

Bei Bedarf können die Kopfzeilen individuell angepasst werden.
Dies geschieht mit den Befehlen \Macro{ihead} und \Macro{ohead}, welche
auf den gleichnamigen Kommandos aus dem Koma-Skript aufbauen und für eine
korrekte vertikale Positionierung sorgen.

\begin{Declaration}
  \Macro{headsepline}
\end{Declaration}


Die standardmäßig gesetzten kurzen Linien am jeweils oberen äußeren Ende des
Kopfbereiches sind im Makro \Macro{headsepline} definiert.
Dies kann bei Bedarf überschrieben werden.


\section{Seitenlayout}

Eine Seite im Corporate Design unterteilt sich grundlegend in 2 Bereiche.
Der Absenderbereich und der Kommunikationsbereich. Der Absenderbereich wird zur
Darstellung des TU-Siegelbandlogos und des Logos eines speziellen Instituts
oder einer Abteilung verwendet. Im Textteil dient er teilweise als Bereich für
Inhalte einer Kopf-/Fußzeile.

Der Absenderbereich kann entweder am Anfang oder am Ende der Seite platziert
werden. Dies beeinflusst auch das Gaußraster, es beginnt jeweils mit dem
größten Segment am Absenderbereich.

\begin{Declaration}
  \Option{a0paper}\\
  \Option{a1paper}\\
  \Option{a2paper}\\
  \Option{a3paper}\\
  \Option{a4paper}\\
  \Option{a5paper}\\
  \Option{a6paper}
\end{Declaration}

Für die Auswahl des verwendeten Papierformats stehen alle DIN A-Größen von
0 bis 6 zur Verfügung. Standardeinstellung ist \Option{a4paper}.

\begin{Declaration}
  \Option{landscape}
\end{Declaration}

Schaltet das Dokument in Querformat-Darstellung.

\begin{minipage}{0.45\textwidth}\sffamily\centering
  {\fboxsep0mm\fbox{\includegraphics[width=\textwidth]{examples/article_landscape.pdf}}}\\
  Dokument im Querformat
\end{minipage}


\begin{important}
  Im Querformat hat das Gaußraster eine abweichende
  Segmentanzahl. Für alle gängigen Formate sind im Querformat \emph{6 Segmente}
  definiert. Daher lösen Layouts, die für hochformatige Dokumente
  geschrieben wurden, im Querformat einen Fehler aus.
\end{important}


\begin{Declaration}
  \KOption{sender}\PName{Position}
\end{Declaration}

Mit der Option \Option{sender} kann die Position des Absenderbereiches
festgelegt werden. Mit der Einstellung \OptionValue{sender}{top} wird
der Absenderbereich am oberen Seitenende dargestellt. Dies ist auch die
Standardeinstellung. Wählt man dagegen \OptionValue{sender}{bottom}, so
wird der Absenderbereich am unteren Ende der Seite platziert

\begin{important}
  Die Positionierung des Absenderbereichs beeinflusst die Orientierung des
  Gaußrasters. Das größte Segment wird dabei immer so gesetzt, dass es direkt
  an den Absenderbereich anschließt und alle folgende von absteigender Höhe sind.
\end{important}



\begin{Declaration}
  \KOption{bcor}\PName{Wert}
\end{Declaration}

Die Bindekorrektur beschreibt einen zusätzlichen Abstand des eigentlichen
Darstellungsbereichs vom inneren Formatrand. Die Bindekorrektur
ist bei Textdokumenten standardmäßig auf Rahmenbreite voreingestellt.

Sinnvoll ist eine Bindekorrektur selbstverständlich zum einen für Bindungen,
wo sie der Breite des durch die Bindung verdeckten Bereiches entsprechen sollte.
Zum anderen ist sie aber auch für den Druckvorgang sinnvoll, da normale Drucker
keinen randlosen Druck ermöglichen. Die Bindekorrektur verhindert so \zB ein 'Abschneiden' des Siegellogos.
Soll das Dokument ohne Bindekorrektur dargestellt werden, so ist dies
mit \OptionValue{bcor}{0mm} möglich.

\begin{Declaration}
  \Option{twoside}
\end{Declaration}

Für das Setzen von zweiseitige Dokumente ist die Option \Option{twoside}
vorgesehen. Mit ihr wird von einseitigem Layout auf zweiseitiges Layout
umgeschaltet, was bedeutet, dass die Innenseite eines Dokuments abwechselnd
auf der linken und rechten Seite definiert ist. Die hat unter anderem Einfluss
auf definierte Ränder (Bindungskorrektur, Marginale).


\begin{Declaration}
  \Option{marginleft}\\
  \Option{marginright}
\end{Declaration}

Das setzen einer Marginale wird durch die Befehle \PName{marginleft} und
\PName{marginright} vereinfacht. Diese setzen jeweils auf der linken (inneren)
bzw. rechten (äußeren) Seite des Dokumentes eine Marginale, deren Breite
einem Element im Spaltenraster entspricht.


\section{Schrift}

Dokumente im Corporate Design werden allgemein in der Schrift \emph{Nexus}
gesetzt.
Diese wird standardmäßig bei allen Dokumenten geladen. Auf eine explizite
Angabe der Schriftgröße in den Dokumentenoptionen sollte nach Möglichkeit
verzichtet werden, da für die jeweiligen Papierformate spezielle
Schriftgrößen vordefiniert sind und geladen werden.
Das Laden anderer Schriftgrößen führt im Vergleich dazu teilweise zu stark abweichenden Ergebnissen.

Für wichtige Standard-Elemente gibt es des weiteren vordefinierte Koma-Fonts.


\begin{desctable}
\entry{\PValue{headline}}{%
  Einfache Überschrift (groß)
}
\entry{\PValue{subheadline}}{%
  Unterüberschrift
}
\entry{\PValue{institute}}{%
  Institutsname im Logo-Bereich des Absenders.
}
% \entry{\PValue{infotext}}{%
% }
% \entry{\PValue{infotextitem}}{%
% }
\entry{\PValue{copytext}}{%
  Mengentext
}
\end{desctable}


\subsection{Farben}

Für eine Übersicht über die verfügbaren Farben siehe \ref{sec:colors}.%TODO...

\begin{Declaration}
  \Option{mono}
\end{Declaration}

Schwarz-Weiß-Darstellung der Standardelemente

\begin{Declaration}
  \Option{cmyk}
\end{Declaration}

CMYK-Darstellung der Standardelemente
   % Kapitel: 'Dokumente'

% \chapter{Briefe}

Briefe können mit der Dokumentenklasse \newdocumentclass{tubslttr2} erstellt
werden.

Der Aufbau eines Briefes unterscheidet sich zum Teil grundlegend von dem eines
normalen Dokumentes, weswegen hierfür eine spezielle Klasse vorliegt%
\footnote{Das etwas spezielle Interface liegt in der zugrunde liegenden 
KOMA-Skript-Klasse scrlttr2 begründet.}, deren Benutzung sich ein wenig von
der der anderen unterscheidet.

Die Absender-Informationen eines Briefes können mit Hilfe einer Reihe von
Variablen festgelegt werden. Näher geht darauf das Kapitel \ref{sec:lettervars}
ein.

Absender-Informationen variieren naturgemäß von Institut zu Institut
und vor allem von Person zu Person. Innerhalb eines Instituts und insbesondere
bei gleichem individuellen Absender sind diese Informationene jedoch sehr
konstant. Daher gibts es die Möglichkeit, die Absender-Informationen in
separaten Dateien als Vorlage vorzudefinieren. Mehr Informationen dazu gibt
das Kapitel \ref{sec:lco}

\section{Paketoptionen}

Generell können die meisten Optionen, die auch die zugrundeliegenden Klasse
\texttt{scrlttr2} akzeptiert, übergeben werden. Es wird jedoch keine korrekte
Funktionalität gewährleistet und darauf hingewiesen, dass diese Optionen
ggf. nich Corporate Design-Konform sind.

Als Schritgröße ist standardmäßig \PValue{10pt} eingestellt.

\begin{Declaration}
  \Option{arial}\\
  \Option{nexus}
\end{Declaration}

Standardmäßig wird ein Brief in der Schriftart \emph{Nexus} gesetzt.
Ist stattdessen eine Darstellung in der Schriftart \emph{Arial} erwünscht,
so kann dies mittels der Option \PName{arial} festgelegt werden.
Die Option \PName{nexus} sorgt für eine Darstellung in \emph{Nexus} und
kann optional angegeben werden.

\begin{Declaration}
  \Option{mono}
\end{Declaration}

Falls ein Brief nur in schwarz-weiß erstellt werden soll, können mit Hilfe
der Option \PName{mono} alle wesentlichen Elemente der Briefvorlage monochrom
dargestellt werden. Dazu zählen unter anderem Das Siegelband-Logo und die
angrenzende Trennlinie.
Auch sämtliche Elemente, die die Farbe \texttt{tuRed} verwenden, werden
in Schwarz statt in Rot dargestellt.

\section{Variablen}\label{sec:lettervars}

Eine Übersicht über alle zu verwendenden Variablen gibt
Tabelle~\ref{table:lettervars}. Es ist dabei meist nicht nötig zu wissen
wo die entsprechenden Werte verwendet werden, da das Design im allgemeinen
festgelegt ist. Zur besseren Veranschaulichung kann jedoch
Abbildung~\ref{fig:lettervarpos} zu Rate gezogen werden.

\begin{desctable}
  \caption{Von der Klasse \newdocumentclass{tubslttr2} unterstützte
    Variablen -- Inhaltsunabhängig}\label{table:lettervars}\\
  \Endfirsthead
  \caption[]{Von der Klasse \newdocumentclass{tubslttr2} unterstützte Variablen
    (\emph{Fortsetzung})}\\
  \Endhead%
  \standardfoot%
  \ventry{fromlogo}{Anweisungen zum Setzen des Institutslogos (Text oder Grafik)}%
  \ventry{frominstitute}{Name des Instituts}%
  \ventry{fromstreest}{Straßennahme und Hausnummer der Absenders}%
  \ventry{fromzipcode}{Postleitzahl des Absenders}%
  \ventry{fromtown}{Absender-Ort}%
  \ventry{fromfaxdirect}{Fax-Durchwahl}%
  \ventry{fromurl}{Internetaddresse}%
  \ventry{frombank}{Bankverbindung des Absenders}%
  \ventry{fromIBAN}{IBAN des Absender}%
  \ventry{fromBIC}{BIC-Code des Absenders}%
  \ventry{fromUStID}{Umsatzsteuer-ID des Absenders}%
  \ventry{fromSteuernummer}{Steuernummer des Absenders}%
  %
  \ventry{fromdepartment}{ggf. Abteilung oder andere Untereinheit des Instituts}%
  \ventry{fromtitle}{Titel des Absenders (Dr., Prof. Dr., \ldots)}
  \ventry{fromname}{Vollständiger Name des Absenders}
  \ventry{signature}{Signatur unter Unterschrift und Grußformel}
  \ventry{fromphonedirect}{Durchwahl}
  \ventry{fromemail}{E-Mail-Adresse}
\end{desctable}

\begin{desctable}
  \caption{Von der Klasse \newdocumentclass{tubslttr2} unterstützte
    Variablen -- Inhaltsabhängig}\label{table:lettervars}\\
  \Endfirsthead
  \ventry{specialmail}{Spezielle Versandart, wird zwischen Rücksendeadresse
    und Empfängeradresse dargestellt.}
  \ventry{yourref}{Feld \glqq Ihr Zeichen:\grqq}
  \ventry{yourmail}{Feld \glqq Ihre Nachricht vom:\grqq}
  \ventry{myref}{Feld \glqq Unser Zeichen:\grqq}
  \ventry{mymail}{Feld \glqq Unsere Nachricht vom:\grqq}
\end{desctable}


\begin{figure}[!ht]
\centering\fboxsep0mm\fbox{%
\includegraphics[width=0.9\textwidth]{../../letter/doc/examples/example_vars.pdf}
}
\caption{Darstellung der Variablen}
\label{fig:lettervarpos}
\end{figure}

\section{lco-Dateien}\label{sec:lco}

Vorschlag: Aufteilung nach Institut, Mitarbeiter.

Aufteilung ist nicht immer eindeutig.

Eine Beispielhafte Aufteilung findet sich im Anhang...

      % Kapitel: 'Brief'

% \chapter{Präsentationen}
\Index{Präsentation}
\Index{Beamer}

% \section{Anwendung und Folienaufbau}

Grundsätzlich handelt es sich bei der Vorlage für Präsentationen lediglich um
ein \emph{theme} für das \emph{beamer}-Paket.
Für allgemeine Fragen zur Präsentationserstellung wir daher auf die
entsprechende Dokumentation\cite{beamer} verwiesen.

Beschrieben werden hier alle Besonderheiten des Corporate-Design-Themes.
Außerdem sollen einige damit verbundene allgemeine Hinweise gegeben werden.

Die Vorlage wird mit dem Beamer-Befehl \lstinline!\usetheme{tubs}! geladen.

\begin{hint}\sloppy
  Alle tubaltex-spezifischen Template-Argumente sollten nicht der Dokumentenklasse,
  sondern bevorzugt beim Laden des entsprechendes Themes übergeben werden,
  um eine klare und wiederverwendbare Struktur zu erhalten.
  
  Zum Beispiel: \lstinline!\usetheme[<Optionen>]{tubs}!
\end{hint}
% TODO: Template argumente von Macro-Argumenten unterscheidbar machen?
  

\section{Titelfolie}\label{sec:titelfolie}

Die Titelfolie ist im Absender-/Kommunikations-Bereich-Layout mit
\gls{glos:siegelbandlogo} im Sinne der allgemeinen Gestaltungsprinzipien des
Corporate Design gehaltens.

Die Kommunikationsfläche ist nach Vorlage des \gls{glos:gaussraster}s
in drei Bereiche aufgeteilt:
Ein \emph{Bildbereich}, der ein Foto oder eine Grafik als Blickfang enthalten
sollte, darunter der \emph{Titelbereich}, der Präsentationstitel,
sowie alle relevanten Informationen trägt
und zum Abschluss ein einfarbiger (roter) Streifen.

Zusätzlich kann in der rechten oberen Ecke des Absenderbereichs ein
Instituts-Logo platziert werden.

\begin{minipage}{0.5\textwidth}
\begin{verbatim}
\title{Corporate Design}
\subtitle{Jetzt mit \LaTeX}
\author{Max Mustermann}

\begin{frame}[plain]
\titlepage
\end{frame}
\end{verbatim}
\end{minipage}
\fboxsep0mm
\begin{minipage}{0.5\textwidth}
\fbox{\includegraphics[width=0.9\textwidth]{examples/titelseite.pdf}}
\end{minipage}

\paragraph{Standardbefehle}

Es können die Standardbefehle zur Titelseitenerstellung, wie
\lstinline{\title},
\lstinline{\subtitle},
\lstinline{\author},
\lstinline{\institute},
\lstinline{\titlegraphic}
und \lstinline{\logo} verwendet werden.

Die Titelseite wird normal mit \lstinline{\titlepage} erzeugt:

\begin{lstlisting}
\begin{frame}[plain]
  \titlepage
\end{frame}
\end{lstlisting}

\emph{Wichtig:} Der \lstinline{frame}-Umgebung muss die Option
\lstinline{[plain]} übergeben werden,
sonst kommt es zu einer falschen Darstellung!

\begin{Declaration}
  \Option{colorfoot}
\end{Declaration}

Mit der Beamer-Option \Option{colorfoot} wird die untere einfarbige Linie
der Titelseite statt in rot in der Farbe des benutzten Sekundärfarbklangs
dargestellt. Weitere Informationen zur Farbanpassung liefert
Kapitel~\ref{beamer:sec:color}.


\subsection{Titelgrafik}
\Index{Titelgrafik!Beamer}
\CommandIndex{titlegraphic}

\begin{Declaration}
  \Macro{titlegraphic}\OParameter{Optionen}\Parameter{Inhalt}
\end{Declaration}

\begin{sloppypar}
Dieser von \LaTeX-beamer bereitgestellte Befehl erlaubt in \tubslatex
die Darstellung von Inhalten im Bildbereich der Titel-Folie.
Der \PName{Inhalt} kann dabei ein beliebiges \LaTeX-Konstrukt sein
und wird am unteren Rande des Bildbereiches ausgerichtet dargestellt.

Im Normalfall wird als Titelgrafik allerdings eine Bild-Datei mittels
\lstinline{\includegraphics} eingebunden werden.
Diese kann (im Rahmen der allgemeinen Corporate Design-Richtlinien)% TODO:cite?
frei gewählt werden.
Das eingebundene Bild wird dabei automatisch randlos
auf den Darstellungsbereich skaliert und zugeschnitten, sofern keine
Standard-Optionen an \lstinline{\includegraphics} übergeben werden oder
manuell eine andere Einpassungsart gewählt wurde.

Innerhalb von \lstinline!\titlegraphic{}! stehen folgende
Optionen zur Einpassung des Bildes zur Verfügung:
\end{sloppypar}

%TODO: float-table + Link?
\begin{desctable}
\toprule
\entry{\PValue{clipped}}{%
  Automatisches Abschneiden. Dies ist die \emph{Standardeinstellung}.
  Das Bild wird dabei optimal in den Darstellungsbereich eingepasst
  (automatische Wahl von \PValue{hclip}/\PValue{vclip}).
}
\entry{\PValue{hclip}}{%
  Das Bild wird vertikal auf Höhe des Bildbereiches skaliert und (falls nötig)
  horizontal zugeschnitten.
}
\entry{\PValue{vclip}}{%
  Das Bild wird horizontal auf Breite des Bildbereiches skaliert und (falls nötig)
  vertikal zugeschnitten.
}
\entry{\PValue{scaled}}{%
  Horizontale \emph{und} vertikale Skalierung. Das Seitenverhältnis des
  Bildes kann dabei verändert und das Bild somit verzerrt werden.
}
\entry{\PValue{keepsize}}{%
  Es wird keinerlei Skalierung oder Beschneidung durchgeführt.
}
\bottomrule
\caption{Mögliche Parameter für \Macro{titlegraphic} zur
automatische Einpassung von Titelgrafik}
\end{desctable}

\begin{Example}
  \begin{lstlisting}
\titlegraphic[scaled]{\includegraphics{myimage.jpg}}
  \end{lstlisting}
  Fügt \textit{myimage.jpg} passend in den Bildbereich ein, indem es das Bild
  horizontal und vertikal skaliert.
\end{Example}

Im Prinzip ist das Seitenverhältnis des Quellbildes durch die automatische
Einpassung nicht besonders kritisch. Jedoch wird durch ein korrektes Verhältnis
sicher gestellt, dass keine wichtigen Teile des Bildes abgeschnitten werden.
In Tabelle \ref{tab:picratio} sind daher Bild-Seiten\-verhältnisse für
verschiedene Seitenverhältnisse der Präsentation aufgeführt.

\begin{table}[ht]
\centering
\begin{tabular}{ll}
\toprule
\bfseries Präsentation  & \bfseries  Bild  \\
\midrule
$3:4$   & $1:3{,}15$ \\
% Seitenverhältnis für 5:4 einfügen
$16:9$  & $1:4{,}30$ \\
$16:10$ & $1:3{,}83$ \\
\bottomrule
\end{tabular}
\caption{Bild-Seitenverhältnisse}
\label{tab:picratio}
\end{table}



\paragraph{Standardgrafik}
\Index{Standardgrafik!Beamer}
\CommandIndex{tuDefaultTitlegraphic}

Zu Testzwecken oder falls kein eigenes Bild zur Hand ist,
kann alternativ ein Standardbild eingefügt werden, das die Front
des TU-Altgebäudes zeigt und gut mit den Standard-Folienfarben harmoniert.

Das Einfügen des Standardbildes ist mit dem Befehl
\linebreak\lstinline{\tuDefaultTitlegraphic} möglich,
der \lstinline!\titlegraphic{}! direkt als Argument übergeben werden kann.

\begin{example}
\begin{lstlisting}
\titlegraphic{\tuDefaultTitlegraphic}
\end{lstlisting}
\end{example}


\paragraph{Manuelle Einpassung}\hfill

\begin{Declaration}
  \Macro{titlegraphicswidth}\\
  \Macro{titlegraphicsheight}
\end{Declaration}

Zur manuellen Einpassung oder für das Erstellen von Grafiken mittels
\LaTeX-Befehlen ist die Dimension des Titelgrafik-Bereichs (Bildbereich) in den
beiden Längen \lstinline{titlegraphicswidth} und \lstinline{titlegraphicsheight}
hinterlegt.

% Das Titelbild wird generell an der unteren Kante des dafür vorgesehenen Bereichs ausgerichtet, sodass in falschem Seitenverhältnis vorliegene Bilder nach Oben 'rausgeschoben' werden.

\subsection{Logo}
\Index{Logo!Beamer}
\Index{Sekundärlogo!Beamer}

Neben dem immer vorhandenen Siegelband-Logo der TU kann noch ein sekundäres Logo
dargestellt werden, welches das jeweilige Institut bzw. die jeweilige Abteilung
repräsentiert.

\begin{Declaration}
  \Macro{logo}\Parameter{Bild}
\end{Declaration}

Das mittels \lstinline{\logo} eingebundene Logo wird in der rechten oberen
Ecke des Absenderbereichs angezeigt und standardmäßig auch auf allen weiteren
Folien im Fußbereich, sofern nicht die Option \Option{nologoinfoot}
verwendet wird.

Das eingebundene Logo wird bei Verwendung von \Macro{includegraphics}
automatisch korrekt vertikal skaliert,
sofern keine manuelle Skalierung als Option angegeben wird.

% Die Verwendung von \lstinline{\logoheight} als Höhenbeschränkung für Grafiken
% stellt deren automatische korrekte vertikale Skalierung sicher.
% Das Seitenverhältnis der verwendeten Grafik-Datei ist dabei relativ
% frei, sollte jedoch nach Möglichkeit und Gründen der Lesbarkeit breiter sein als
% hoch.

\begin{example}
\begin{lstlisting}
\logo{\includegraphics{institut.jpg}}
\end{lstlisting}
\end{example}

\begin{Declaration}
  \Macro{logoheight}
\end{Declaration}

Für allgemeine Inhalte kann auf die Höhe des Logo-Bereichs über die Länge
\Macro{logoheight} zugegriffen werden.

\section{Inhaltsfolien}
\Index{Folien!Inhalt}
\Index{Inhaltsfolien}

Das Layout der Inhaltsfolien ist relativ schlicht und bietet viel Platz.
Es zeichnet sich durch einen Kopfbereich, der Folientitel und evtl. eine
Inhaltsübersicht enthält, und einen Fußbereich, der neben Siegelband- und
Institutslogo allgemeine Informationen zur Präsentation liefert.
Auf den genaueren Aufbau des Kopfbereichs geht Kapitel~\ref{subsec:head},
auf den des Fußbereichs Kapitel~\ref{subsec:foot} ein.
% Die Überschrift ist im Kopfbereich grau hinterlegt und im Fußbereich der Folien
% findet sich das TU-Logo, sowie eine rote Trennlinie.
% Rechts unterhalb der Trennlinie wird, wenn definiert, das Logo angezeigt.
% Ist dies nicht gewünscht, kann dies mit der Klassenoption
% \Option{nologoinfoot} unterdrückt werden.

\begin{minipage}{0.5\textwidth}
\begin{verbatim}
\begin{frame}{Inhaltsseite}
  \begin{itemize}
    \item Hier steht der Inhalt
    \item Hier nicht
    \item Weitere Informationen
  \end{itemize}
\end{frame}
\end{verbatim}
\end{minipage}
\begin{minipage}{0.5\textwidth}
\fboxsep0mm
\fbox{\includegraphics[width=0.9\textwidth,page=1]{examples/inhaltsseite.pdf}}
\end{minipage}

\subsubsection{Hervorgehobene Folien}

Einzelne wichtige Folien können extra hervorgehoben werden.
Dazu dient die Option \Option{highlight},
welche einen breiteren und rot hinterlegten Titelbereich erzeugt.

\begin{minipage}{0.5\textwidth}
\begin{verbatim}
\begin{frame}[highlight]
  {Inhaltsseite -- Hervorgehoben}
  \begin{itemize}
    \item Hier steht der Inhalt
    \item Hier nicht
    \item Weitere Informationen
  \end{itemize}
\end{frame}
\end{verbatim}
\end{minipage}
\begin{minipage}{0.5\textwidth}
\fboxsep0mm
\fbox{\includegraphics[width=0.9\textwidth,page=2]{%
  examples/inhaltsseite.pdf}}
\end{minipage}

\subsubsection{Titellose Folien}

Wird bei einer Folie kein Titel angegeben, so wird der graue Titelbereich
nicht dargestellt.
Ein eventuell vorhandenes Inhaltsverzeichnis wird jedoch weiterhin samt
Hintergrund angezeigt!

% TODO: Option notitle?
\begin{minipage}{0.5\textwidth}
\begin{verbatim}
\begin{frame}
  \begin{itemize}
    \item Hier steht der Inhalt
    \item Hier auch
    \item Weitere Informationen
  \end{itemize}
\end{frame}
\end{verbatim}
\end{minipage}
\begin{minipage}{0.5\textwidth}
\fboxsep0mm
\fbox{\includegraphics[width=0.9\textwidth,page=3]{%
  examples/inhaltsseite.pdf}}
\end{minipage}

\subsection{Kopfbereich}\label{subsec:head}
\Index{Kopfbereich!Beamer}

Der Kopfbereich der Folien besteht standardmäßig aus dem grau hinterlegten
Folientitel. Optional kann darüber auch noch eine Inhaltsübersicht
(Inhaltsverzeichnis) dargestellt werden, 
die auch zur Navigation im Dokument genutzt werden kann.
Dieser Kopfbereich ist dann unabhängig vom Titelbereich und wird
auch bei titellosen Folien dargestellt.

\begin{Declaration}
  \Option{colorhead}
\end{Declaration}

Der Kopfbereich ist normalerweise in einem hellen Grauton hinterlegt.
Durch Verwendung der Option \Option{colorhead} wird der Hintergrund
in der 40\%-Version der aktiven Variante des Sekundärfarbklangs gesetzt.
Nähere Informationen und Beispiele zur Farbanpassung liefert das
Kapitel~\ref{beamer:sec:color}.

\subsubsection{Inhaltsübersicht}

Die Inhaltsübersicht (Inhaltsverzeichnis) wird bei Verwendung rechts
ausgerichtet als oberstes Element jeder Folie dargestellt.
Sie wird auch auf titellosen Folien angezeigt.


\begin{Declaration}
  \Option{tocinheader}\\
  \Option{tinytocinheader}
\end{Declaration}

Das Inhaltsverzeichnis im Folienkopf
wird mir der Option \Option{tocinheader} aktiviert.
Die \Option{tinytocinheader} bewirkt dies ebenfalls,
verwendet aber eine kleinere Schriftgröße.

Angezeigt werden standardmäßig untereinander die Gliederungsebenen
\Macro{section} und \Macro{subsection}. Dies kann mit der Option
\Option{nosubsectionsintoc} deaktiviert werden.


\begin{Declaration}
  \Option{widetoc}\\
  \Option{narrowtoc}
\end{Declaration}

Während die Option \Option{widetoc} das Inhaltsverzeichnis in der Breite
streckt, wird es durch die Option \Option{narrowtoc} gestaucht.
Somit kann eine flexiblere Platzausnutzung erreicht werden.

\begin{Declaration}
  \Option{nosubsectionsintoc}
\end{Declaration}

Soll im Inhaltsverzeichnis des Kopfbereichs die Gliederungsebene
\Macro{subsection} nicht mit angezeigt werden, so kann dies durch die Option
\Option{nosubsectionsintoc} erreicht werden.

Dies Option ist natürlich nur wirksam, wenn die Option
\lstinline{tocinheader} oder \lstinline{tinytocinheader} verwendet wird.



\subsection{Fußbereich}\label{subsec:foot}
\Index{Fußbereich!Beamer}

Der Fußbereich zeichnet sich durch das Siegelband-Logo und eine rote Trennlinie
aus. Darunter stehen standardmäßig Informationen wie Datum, Autor(en), Präsentationstitel
und Seitenzahl, sowie auf der rechten Seite ein Instituts-/bzw. Abteilungslogo,
sofern definiert.
Der Inhalt der einzelnen Felder ergibt sich aus der Benutzung der in
\chaptername~\ref{sec:titelfolie} beschriebenen Befehle.
Die konkrete Darstellung in der Fußzeile kann jedoch durch eine Reihe von
Beamer-Optionen zusätzlich angepasst werden.

Autor und Titel werden, falls vorhanden, in ihrer Kurzversion dargestellt,
welche jeweils bei \Macro{author} und \Macro{title} im optionalen Argument
definiert werden kann.

Die Darstellung der einzelnen Textelemente wird abhängig von deren
jeweilige Länge durch ein dynamisches Positionierungsschema optimal
an den verfügbaren Platz angepasst.
%TODO: Details zu automat. Umbrüchen etc.?

% Sprache
Ändert man im Babel-Paket die Sprache, ändert sich auch automatisch
die Sprache der Fußzeile (z.B. "`page"' statt "`Seite"', anderes Datumsformat, etc.).
Vollständig unterstützt sind Deutsch und Englisch.

\begin{Declaration}
  \Option{nodate}\\
  \Option{noauthor}\\
  \Option{nopagenum}
\end{Declaration}

Diese Beamer-Optionen unterdrücken jeweils die Darstellung eines Elementes in der Fußzeile.
Die Option \Option{nodate} sorgt dafür, dass das Datum nicht dargestellt wird,
\Option{noauthor} dafür, dass der Autor nicht angezeigt und
\Option{napagenum} dafür, dass die Seitenzahl ausgeblendet wird.

\begin{Declaration}
  \Option{totalpages}
\end{Declaration}

Standardmäßig wird die Seitenzahl wie folgt dargestellt: "`Seite x"'.
Die Option \Option{totalpages} bewirkt, dass zusätzlich die Gesamtseitenzahl
mit eingeblendet wird: "`Seite x von y"'.

\begin{Declaration}
  \Option{nologoinfoot}
\end{Declaration}

Diese Beamer-Option deaktiviert die Darstellung des Institutslogos im Fußbereich.

\subsection{Inhalt}

\paragraph{Blöcke}

Blöcke werden in den Vorlagen bei Standardeinstellungen einfach
mit transparentem Hintergrund mit verschiedenfarbigen Überschriften dargestellt
wie in \figurename~\ref{subfig:stdblocks} zu sehen.
Die Farben sind dabei dem CD-Farbschema entnommen.

\begin{Declaration}
  \Option{colorblocks}
\end{Declaration}

Für eine farblich abgesetztere Darstellung steht die Option \Option{colorblocks}
zur Verfügung. Diese bewirkt, dass die Blöcke in farbigen Boxen dargestellt werden
(siehe \figurename~\ref{subfig:colorblocks}).


\begin{figure}[!ht]
\begin{minipage}{0.5\textwidth}
  \centering
  \fboxsep0mm\fbox{%
  \includegraphics[width=0.95\textwidth,page=1]{examples/beamer_blocks.pdf}}
  \subcaption{Standard-Blöcke}\label{subfig:stdblocks}
\end{minipage}
\begin{minipage}{0.5\textwidth}
  \centering
  \fboxsep0mm\fbox{%
  \includegraphics[width=0.95\textwidth,page=2]{examples/beamer_blocks.pdf}}
  \subcaption{\Option{colorblocks}-Blöcke}\label{subfig:colorblocks}
\end{minipage}
\caption{Vergleich zwischen normaler und \Option{colorblocks}-Darstellung}
\end{figure}


\paragraph{Beschriftung von Abbildungen}
\Index{Abbildungen!Beamer}
\Index{caption!Beamer}

Anders als bei den Standard-Templates werden Abbildungen
automatisch durchnummeriert.

Sollte dies nicht erwünscht sein, kann dies mittels
\begin{lstlisting}
\setbeamertemplate{caption}[default]
\end{lstlisting}
deaktiviert werden.

\section{Farbanpassung}\label{beamer:sec:color}
\Index{Farben!Beamer}

Um ein individuelleres Foliendesign zu ermöglichen, können einige Elemente
(unter Verwendung des CD-Farbkontingents)
farblich angepasst werden.
Dazu zählen die Elemente der Titelseite (mit Ausnahme des Logos),
sowie die Hintergrundfarbe der Folientitel.

Zur einfachen Farbanpassung stehen einige spezielle Optionen zur Verfügung.
Eine komplett freie Gestaltung kann bei Bedarf über Änderung der
jeweiligen Templates erfolgen. Beide Möglichkeiten werden im folgenden
erläutert.

Details über die Farben der CD-Farbklänge können dem
Kapitel~\ref{chap:tubscolors} entnommen werden.

\subsection{Anpassung über Optionen}

Über die Angabe von Farbargumenten können alle zur farblichen Anpassung
vorgesehenen Elemente nach einem vordefinierten \emph{Farbschema}
verändert werden.
Dazu kann ein Farbklang sowie eine Variante gewählt werden.

Die Optionen gliedern sich in 3 Kategorien:
\begin{itemize}
  \item \emph{Farbe} -- Auswahl des Farbklangs\\
    Optionen: \texttt{yellow}, \texttt{green}, \texttt{blue}, \texttt{violet}
  \item \emph{Helligkeit} -- Auswahl der Helligkeitsabstufung
    innerhalb des gewählten Farbklangs\\
    Optionen: \texttt{dark}, \texttt{medium}, \texttt{light}
  \item \emph{Elemente} -- Färbung von Fußbereich/Kopfbereich aktivieren\\
    Diese wurden bereits in \chaptername~\ref{sec:titelfolie} und \ref{subsec:head} beschrieben.\\
    Optionen: \texttt{colorhead}, \texttt{colorfoot}
\end{itemize}

% Weitere Details zu Verwendung der Optionen sind im Kapitel \ref{beamer:optionen}
% zu finden.

\paragraph{Farbklang}
Es stehen vier verschiedene Farbklänge mit aufeinander abgestimmten Farbtonwerten
und Varianten zur Verfügung.

\begin{Declaration}
  \Option{orange}~/~\Option{yellow}\\
  \Option{green}\\
  \Option{blue}\\
  \Option{violet}
\end{Declaration}

Mit den Optionen \Option{orange} (alternativ \Option{yellow}), \Option{green}, \Option{blue}
oder \Option{violet} erfolgt die Auswahl des
gelben, grünen, blauen oder violetten Farbklangs als Farbschema.

\paragraph{Variante}
Für jeden Farbklang gibt es 3 verschiedene Varianten, die sich in ihrem Ton
und ihrer Helligkeit unterscheiden.

\begin{Declaration}
  \Option{dark}\\
  \Option{medium}\\
  \Option{light}
\end{Declaration}

Auswahl der hellen, mittleren bzw. dunklen Farbreihe aus gewähltem Farbklang.

\begin{example}
\lstinline!\usetheme[green,light]{tubs}! Wählt den grünen Farbklang und die
helle Farbreihe. Damit wird das Farbschema auf hellgrün gesetzt.
\end{example}

\paragraph{Farbmodell}
Standardmäßig werden Folien in einem optimierten RGB-Farbmodell ausgegeben.
Bei Bedarf kann ein anderes Farbmodell zur Darstellung gewäht werden.


\begin{Declaration}
  \Option{rgb}\\
  \Option{rgbprint}\\
  \Option{cmyk}
\end{Declaration}

Mit der Option \Option{rgb} erfolgt die Ausgabe im Beamer-optimierten RGB-Farbmodell.
Dies entspricht der Standardeinstellung.
Mit der Option \Option{rgbprint} erfolgt die Ausgabe in RGB-Druckfarben.
Mit der Option \Option{cmyk} erfolgt die Ausgabe im CMYK-Farbmodell.

\clearpage
\subsubsection{Beispiele}

% \begin{minipage}{0.55\textwidth}
% \lstinline!\usetheme[green,dark]{tubs}!
% \end{minipage}\hfill
\begin{center}
\begin{minipage}{0.49\textwidth}
\fboxsep0mm
\fbox{\includegraphics[width=0.99\textwidth,page=1]{%
  examples/colorscheme1.pdf}}
\end{minipage}\hfill
\begin{minipage}{0.49\textwidth}
\fboxsep0mm
\fbox{\includegraphics[width=0.99\textwidth,page=2]{%
  examples/colorscheme1.pdf}}
\end{minipage}\medskip\\
{\ttfamily \textbackslash usetheme[green,dark]\{tubs\}}
\end{center}

\begin{center}
\begin{minipage}{0.49\textwidth}
\fboxsep0mm
\fbox{\includegraphics[width=0.99\textwidth,page=1]{%
  examples/colorscheme2.pdf}}
\end{minipage}
\begin{minipage}{0.49\textwidth}
\fboxsep0mm
\fbox{\includegraphics[width=0.99\textwidth,page=2]{%
  examples/colorscheme2.pdf}}
\end{minipage}\medskip\\
{\ttfamily \textbackslash usetheme[green,dark,colorfoot,colorhead]\{tubs\}}
\end{center}

\begin{center}
\begin{minipage}{0.49\textwidth}
\fboxsep0mm
\fbox{\includegraphics[width=0.99\textwidth,page=1]{%
  examples/colorscheme3.pdf}}
\end{minipage}
\begin{minipage}{0.49\textwidth}
\fboxsep0mm
\fbox{\includegraphics[width=0.99\textwidth,page=2]{%
  examples/colorscheme3.pdf}}
\end{minipage}\medskip\\
{\ttfamily%
  \textbackslash usetheme[blue,colorhead]\{tubs\}}
\end{center}

\clearpage
\subsection{Erweiterte Anpassung über Farb-Templates}

Eine individuellere Farbanpassung ist über die Änderung der entsprechenden
Templates möglich.

\subsubsection{Paletten}

Das Grundfarbschema eines Beamer-Templates wird über sog. Paletten definiert.
Es existieren eine primäre, sekundäre, tertiäre und quaternäre Palette,
die jeweils eine bestimmte Farbgebung beeinflussen.

In den \tubslatex-Vorlagen erfüllen die Paletten die folgenden Funktionen:

\begin{description}
  \item[\ttfamily palette primary]
    Die primäre Palette definiert die Farbgebung
    von Kopfleiste und Titelleiste auf Inhaltsfoline,
    sowie von Part-Seiten.
  \item[\ttfamily palette secondary]
    Die sekundäre Palette definiert die Farbgebung der einzelnen
    Segmente der Titelseite.
  \item[\ttfamily palette tertiary]
    Die tertiäre Palette definiert die Farbe der Fußzeilen-Schriftelemente.
  \item[\ttfamily palette quaternary]
    Die quaternäre Palette wird in den Vorlagen nicht verwendet.
\end{description}

Die Paletten können jeweils mittels \Macro{setbeamercolor}
unter Angabe der Palettenbezeichnung modifiziert werden.

\begin{example}
\begin{lstlisting}
  \setbeamercolor{palette primary}{fg=tuOrange, bg=tuGreen}
\end{lstlisting}
\end{example}

\subsubsection{Einzelelemente}

Neben den Paletten können natürlich auch einzeln Elemente manuell angepasst
werden. Die wichtigsten sind in der nachfolgenden Tabelle aufgeführt:

\begin{tabularx}{\textwidth}{lX}
  Bezeichner & Wirkung\\
  \midrule
  \ttfamily titlebarfirst &
    Der Hintergrund des Grafik-Segments auf der Titelseite\\
  \ttfamily titlebarsecond &
    Das Segment unterhalb des Grafik-Segments auf der Titelseite\\
  \ttfamily titlebarlow &
    Die schmale Fußleiste auf der Titelseite\\
  \ttfamily titlehead &
    Die Kopfzeile der Inhaltsseite\\
  \ttfamily titlelike &
    Die Titelzeile der Inhaltsseite\\
  \ttfamily block body &
    Hintergrundfarbe von Blöcken\\
  \ttfamily block title &
    Titelfarbe von Blöcken\\
  \ttfamily block body alerted &
    Hintergrundfarbe von alerted-Blöcken\\
  \ttfamily block title alerted &
    Titelfarbe von alerted-Blöcken\\
  \ttfamily block body example &
    Hintergrundfarbe von example-Blöcken\\
  \ttfamily block title example &
    Titelfarbe von example-Blöcken\
\end{tabularx}

Die Anpassung weiterer Elemente ist natürlich möglich,
wird aber nicht empfohlen.
Sie können bei Bedarf den entsprechenden Paket-Dokumentationen entnommen werden.


% \section{Template-Optionen}\label{beamer:optionen}
% 
% Auflistung aller template-spezifischer Optionen zur Übergabe an die
% Dokumentenklasse bzw. das Beamer-theme.
% 
% Die template-spezifischen Argumente sollten bevorzugt
% über \lstinline!\usetheme! übergeben werden.
%TODO: Hinweis zu Beamer-Option!?

\section{Schrift}

Die voreingestellte Standardschrift der Vorlagen ist Arial in Schriftgröße 11pt.
Sie kann bei Bedarf auf Nexus gewechselt werden.

\begin{Declaration}
  \Option{arial}\\
  \Option{nexus}
\end{Declaration}

Mit der Option \Option{arial} wird die Schriftart Arial verwendet.
Die Option \Option{nexus} wählt Nexus als Schriftart aus.

\subsection{Zusätzliche Schriftgrößen}
\Index{Schriftgrößen!Beamer}
\CommandIndex{microsize}
\CommandIndex{nanosize}

Das font-Template legt zwei zusätzliche Schriftgrößen \lstinline{\microsize}
und \lstinline{\nanosize} an, die jeweils noch kleiner sind als alle restlich
verfügbaren.

Sie werden intern verwendet, können aber auch bei begründeter
Notwendigkeit von Autoren verwendet werden,
z.B. für kleinste Quellenangaben etc.

\subsection{Elemente anpassen}

Mit Hilfe der Beamer-Fonts können bei Notwendigkeit auch einzelne Elemente
separat in ihrer Schrift angepasst werden.
Dieses Mittel sollte jedoch mit Bedacht eingesetzt werden.
Informationen zur Anpassung können der Dokumentation zum Beamer-Paket entnommen
werden. Die Beschreibung modifizierbarer Elemente findet sich in der
Paket-Dokumentation zum Font-Template.

\section{Hinweise}

\subsection{Inner-, Outer-, Font-, Color-Template}

Die einzelnen Teil-Templates des beamer-Templates können,
entsprechend der beamer-Konventionen, selbstverständlich
auch unabhängig voneinander verwendet werden.
So kann zum Beispiel das color-Template in Kombination mit einem der
Standard-Template verwendet werden, etc.

Die Kombination mit anderen Templates ist im allgemeinen jedoch nicht
CorporateDesign-konform, weswegen davon \emph{ausdrücklich} abgeraten wird.

\subsection{Skalierbarkeit (Format)}
\Index{16:9}
\Index{Format!Beamer}

Die Beamer-Option \lstinline{aspectratio}, mit der die Präsentation auch auf
andere Seitenformate gebrachte werden kann, wird vom Layout generell
unterstützt, jedoch ist zu bedenken, dass dafür ein Titelbild in einem anderen
Seitenverhältnis benötigt wird, wenn es verlustfrei dargestellt werden soll
(siehe Tabelle \ref{tab:picratio}).

\begin{hint}
Die Option \lstinline{aspectratio} steht erst ab beamer-Version 3.10
zur Verfügung.
\end{hint}


\subsection{columns-Umgebung}
\Index{columns-Umgebung}

Die Verwendung der columns-Umgebung ist in vielen Fällen sinnvoll.
Jedoch sollte die Umgebung nur mit der Option \lstinline{onlytextwidth}
verwendet werden, da sonst ein zusätzlicher Abstand zwischen den
Spalten hinzugefügt werden und die Gesamtbreite zu groß wird:

\begin{lstlisting}[morekeywords={onlytextwidth},keywordstyle=\color{tuOrange}]
\begin{columns}[onlytextwidth]
  \column{0.5\textwidth}
  \column{0.5\textwidth}
\end{columns}
\end{lstlisting}


\subsection{Listings in Folien}
\Index{Listings!Beamer}

Auch wenn dies in der beamer-Dokumentation ausdrücklich erläutert ist,
sollte hier noch einmal erwähnt werden, dass bei Gebraucht von Listings etc.
in Folien die Option \lstinline{fragile} übergeben werden muss.

\begin{lstlisting}[morekeywords={fragile},keywordstyle=\color{tuOrange}]
\begin{frame}[fragile]{Titel}
  % Listings, etc.
\end{frame}
\end{lstlisting}



\subsection{PDF-Titel}

Die standardmäßige Erstellung des PDF-Titels von beamer ist dahingehend
abgeändert, dass, sofern vorhanden, die \emph{Kurzversion} des Präsentationstitels
als PDF-Titel verwendet wird.

Die Verwendung der Kurzversion ermöglicht es zum Beispiel,
kleine Grafiken in den Titel einzubringen und
trotzdem einen korrekten PDF-Titel zu erhalten.


\section{Minimalbeispiel}% TODO: Anhang?!

Im folgenden ist der Code eines Minimalbeispiels aufgeführt, zusammen mit den
daraus erzeugten Folien.

\begin{verbatim}
\documentclass{beamer}

\usetheme{tubs}
\usepackage[T1]{fontenc}
\usepackage[utf8x]{inputenc}

\title{Corporate Design}
\subtitle{Jetzt mit \LaTeX}
\author{Max Mustermann}
\titlegraphic{\includegraphics[width=\titlegraphicswidth]{titlepicture}}
\logo{\includegraphics[height=\logoheight]{institut.jpg}}

\begin{document}

\begin{frame}[plain]
  \titlepage
\end{frame}

\begin{frame}{Inhaltsseite}
  \begin{itemize}
    \item Hier steht der Inhalt
    \item Hier nicht
    \item Weitere Informationen
  \end{itemize}
\end{frame}

\end{document}
\end{verbatim}

\begin{center}
  \fbox{\includegraphics[width=0.9\textwidth]{examples/titelseite.pdf}}

  \fbox{\includegraphics[width=0.9\textwidth]{examples/inhaltsseite.pdf}}
\end{center}
        % Kapitel: 'Präsentationen'

% \chapter{Plakate}

Plakate können in \tubslatex mit Hilfe der Dokumentenklasse \newdocumentclass{tubsposter} erstellt werden.
Dabei werden drei unterschiedliche Plakat-Typen unterstützt,
zwischen denen mit der Klassenoption \KOption{style}\PName{<standard/scifi/bulletin>} gewählt werden kann.

Dies sind zum einen \emph{Veranstaltungsplakate},
welche sich am normalen Gaußraster
orientieren und als Standard-Darstellungstyp voreingestellt sind.

Für die Darstellung \emph{wissenschaftlicher} Inhalte gibt es ein spezielles Layout,
das mehr Platz und Flexibilität für die Unterbringung von Informationen bietet.
Es orientiert sich an einem eigenen Modulsystem und nicht am Gaußraster.

Zusätzlich gibt es noch eine Unterstützung für die Erstellung einfacher
\emph{Aushänge}. Diese weisen nur einen einfachen großen Textbereich auf, der
nach oben durch das Siegelband-Logo mit einer Trennlinie begrenzt ist.

\begin{figure}[htp]\centering
\fboxsep0mm
\begin{minipage}{0.32\textwidth}\centering
  \fbox{\includegraphics[width=0.9\textwidth]{examples/posterstyle_standard.pdf}}
  \subcaption{Veranstaltung (\PValue{standard})}
\end{minipage}
\begin{minipage}{0.32\textwidth}\centering
  \fbox{\includegraphics[width=0.9\textwidth]{examples/posterstyle_scifi.pdf}}
  \subcaption{wissenschaftlich (\PValue{scifi})}
\end{minipage}
\begin{minipage}{0.32\textwidth}\centering
  \fbox{\includegraphics[width=0.9\textwidth]{examples/posterstyle_bulletin.pdf}}
  \subcaption{Aushang (\PValue{bulletin})}
\end{minipage}
\caption{Die verschiedenen Plakat-Typen}
\end{figure}

\begin{Declaration}
  \KOption{style}\PName{<standard/scifi/bulletin>}
\end{Declaration}

Mit \KOption{style}{standard} wird die Darstellung für Veranstaltungsplakate
gewählt. Diese ist bereits voreingestellt und muss daher nicht explizit gewählt
werden.
Mit \KOption{style}{scifi} wir das Layout-System für wissenschaftliche Plakate
geladen, welches sich teilweise grundlegend von dem normalen Layout unterscheidet.
Der Wert \KOption{style}{bulletin} erlaubt die schnelle Erstellung von
Aushängen und lädt eine relative schlichte voreingestellte Darstellung.

\begin{Declaration}
  \KOption{sender}\PName{top/bottom}
\end{Declaration}

Diese Option legt fest, ob der Absenderbereich am oberen oder unteren Blattrand
platziert werden soll. Für weitere Informationen siehe vorausgehende Kapitel.
%TODO: ja? wo?
%TODO: Klasse + Umgebung


\begin{Declaration}
  \Option{a3paper}\\
  \Option{a2paper}\\
  \Option{a1paper}\\
  \Option{a0paper}
\end{Declaration}

Für alle Poster-Varianten stehen vier Papierformate zur Auswahl für die
Schriftgrößen etc. vordefiniert sind und automatisch geladen werden.
Dies sind A3, A2, A1 und A0, die mit den entsprechenden Optionen
\Option{a3paper}, \Option{a2paper}, \Option{a1paper} bzw. \Option{a0paper}
gewählt werden können.
\bigskip

Im folgenden werden die einzelnen Poster-Typen einzeln vorgestellt und beschrieben.

\paragraph{Weitere Klassenoptionen:}

Wie bei allen anderen \tubslatex-Klassen können auch für Poster
die diversen verfügbaren Optionen zur Farbdarstellung oder Schriftartwahl
verwendet werden.
Näheres hierzu findet sich in den vorausgehenden Kapiteln oder detaillierter
in Kapitel~\ref{chap:tubscolors} und \ref{chap:nexus}.

\clearpage
\section{Veranstaltungsplakate}

Veranstaltungsplakate werden im Gaußraster gesetzt und entsprechen somit
im Grunde allen Standarddarstellungen des Corporate Design.

Veranstaltungsposter werden entweder ohne Option oder
unter Angabe der Optionswahl \KOption{style}\PName{standard} erstellt:
\begin{lstlisting}
\documentclass[style=standard]{tubsposter}
\end{lstlisting}

Da das Gauß-Layoutsystem zum einen umfangreich, zum anderen auch für verschiedene Zwecke eingesetzt wird und werden kann, ist ihm ein Extra-Kapitel (\ref{chap:gausspage})
gewidmet.
Dies beschreibt die möglichen Befehle und Optionen ausführlich.
Hier werden sie daher jeweils nur kurz angerissen.

% TODO: landscape?

\begin{Declaration}
  \XMacro{begin}\PParameter{\Environment{tubsposter}}%
    \OParameter{options}\\
  \quad\dots\\
  \XMacro{end}\PParameter{tubsposter}
\end{Declaration}

Ein neues Plakat wird mit der Umgebung \Environment{tubsposter} erstellt.
Der optionale Parameter \PName{options} akzeptiert dabei die
unter \ref{sec:gausspage:bglayout} beschriebenen Optionen.

Der Inhalt des Posters sollte aus einem oder mehreren Segmenten bestehen.
Diese können jeweils aus einem oder mehreren zusammengefassten Basis-Segmenten
des Gaußrasters bestehen. Dabei sollten alle Basis-Segmente verwendet werden.

\begin{sloppypar}
Das Siegelband-Logo kann durch Verwendung des Befehls \Macro{showtubslogo}
innerhalb der \Environment{tubsposter}-Umgebung angezeigt werden.
Das Institutslogo wird analog mit \Macro{showlogo}\Parameter{Inhalt}
dargestellt. Hierzu gibt das Kapitel~\ref{subsec:gausspage:bgelement}
weitere Details.
\end{sloppypar}
%TODO: ref, link, ?
% sender
% bgcolor?

\begin{Declaration}
  \XMacro{begin}\PParameter{\Environment{posterrow}}%
    \OParameter{options}
    \Parameter{Höhe}\\
  \quad\dots\\
  \XMacro{end}\PParameter{posterrow}%\\
%   \XMacro{begin}\PParameter{\Environment{posterrow*}}%
%     \OParameter{cols}\\%
%   \quad\dots\\
%   \XMacro{end}\PParameter{posterrow*}
\end{Declaration}

Einzelne Poster-Segmente im Gaußraster können mit der Umgebung
\Environment{posterrow} erzeugt werden.
Dabei gibt der Parameter \PName{Höhe} die Höhe des
Poster-Segments in Basis-Segmenten des Gaußrasters an.
Der optionale Parameter \PName{options} akzeptiert die
unter \ref{subsec:gausspage:bgelement} beschriebenen Optionen mit
denen zum Beispiel recht einfach eine Hintergrundfarbe oder ein Hintergrundbild
gesetzt werden kann.

\begin{figure}[!ht]
\begin{minipage}{0.65\textwidth}
\begin{lstlisting}
\documentclass{tubsposter}
\begin{document}
\begin{tubsposter}
\showtubslogo
\begin{posterrow}[bgcolor=tuGreen]{3}
\usekomafont{headline}Der Titel
\end{posterrow}
\begin{posterrow}[bgcolor=tuOrange]{4}
\end{posterrow}
\begin{posterrow}[bgcolor=tuGreenDark]{1}
\end{posterrow}
\end{tubsposter}
\end{document}
\end{lstlisting}
\end{minipage}
\begin{minipage}{0.35\textwidth}
  \fboxsep0mm
  \fbox{\includegraphics[width=0.98\textwidth]{examples/stdposter_minimal.pdf}}
\end{minipage}
\caption{Minimal-Beispiel zur Verwendung von \Environment{tubsposter} und
  \Environment{posterrow}}
\end{figure}


\clearpage
\section{Wissenschaftliche Plakate}

Wissenschaftliche Poster unterscheiden sich im Layout in diversen Punkten
von allen übrigen Vorlagen. Die Darstellung ist dazu optimiert möglichst
viel Inhalt möglichst flexibel platzieren zu können.

Das Grundlayout weist einen schmaleren Absenderbereich mit einem auf 90\%
skalierten Siegelband-Logo auf.
Der Inhaltsbereich ist rahmenlos und komplett einfarbig hinterlegt.

Der Inhalt kann in sogenannten \emph{Modulen} platziert werden.
Dies sind Kästchen, die durch einen schmalen Abstand optisch
voneinander getrennt sind und so die Informationseinheiten klar voneinander
trennen.
Die Platzierung der Module erfolgt in einem lockeren Rastersystem,
dass die einzelnen Bereiche in Zeilen, Spalten und Unterzeilen teilen kann.

Wissenschaftliche Poster werden unter Angabe der Optionswahl
\KOption{style}\PName{scifi} erstellt:
\begin{lstlisting}
\documentclass[style=scifi]{tubsposter}
\end{lstlisting}

\begin{Declaration}
  \XMacro{begin}\PParameter{\Environment{tubsposter}}%
    \OParameter{options}%
    \Parameter{rows}\\
  \quad\dots\\
  \XMacro{end}\PParameter{tubsposter}
\end{Declaration}

Für die Erstellung von wissenschaftlichen Plakaten wird die Umgebung
\Environment{tubsposter} verwendet, welche
(anders als bei Veranstaltungsplakaten) jedoch einen
zusätzlichen Parameter \PName{rows} erwartet.
Mit diesem wird die Anzahl an Modulzeilen bestimmt, in die der Darstellungsbereich geteilt werden soll.Dies geschieht mittels einer kommagetrennten Liste,
wobei jedes Element entweder eine Länge oder der Buchstabe 'X' sein kann.
Eine Länge legt die Höhe der jeweiligen Modulzeile genau fest,
ein X sorgt dafür, dass alle mit X gekennzeichneten
Zeilen den restlichen zur Verfügung stehenden Platz gleichmäßig untereinander
aufteilen.
Dieses Vorgehen ist an die Tabellen-Umgebung \Environment{tabularx} angelehnt.

\newcommand{\examplegraphic}[2][]{%
  \fboxsep0mm\fbox{\includegraphics[#1]{#2}}
}

\begin{minipage}{0.6\textwidth}
\begin{Example}
  \noindent\Macro{begin}\PParameter{\Environment{tubsposter}}
    \Parameter{6cm,X,4cm,X}\par
  \noindent Erzeugt 4 Modulzeilen, wobei die erste 6cm und die dritte 4cm hoch sind.
  Die Zeilen 2 und 4 nehmen den restlichen verfügbaren Platz ein
  und sind gleich hoch.
\end{Example}
\end{minipage}
\hfill
\begin{minipage}{0.3\textwidth}\centering
  \examplegraphic[width=\textwidth,page=1]{examples/module_split.pdf}
\end{minipage}



\begin{Declaration}
  \XMacro{begin}\PParameter{\Environment{posterrow}}%
    \OParameter{cols}\\%
  \quad\dots\\
  \XMacro{end}\PParameter{posterrow}\\
  \XMacro{begin}\PParameter{\Environment{posterrow*}}%
    \OParameter{cols}\\%
  \quad\dots\\
  \XMacro{end}\PParameter{posterrow*}
\end{Declaration}

Die mit \Environment{tubsposter} angelegten Modulzeilen können nun jeweils mit
der Umgebung \Environment{posterrow} mit Inhalt gefüllt werden.
Dabei kann entweder direkt der gewünschte Inhalt geschrieben oder der optionale
Parameter \PName{cols} benutzt werden.
Dieses erlaubt die Definition zusätzlicher Spalten in der aktuellen Modulzeile.
Es wird wieder ein kommagetrennte Liste erwartet,
deren Elemente dieselbe Bedeutung haben wie bereits beschrieben, außer, dass
sie die Breite und nicht die Höhe definieren.

Der normale Abstand des Inhalts vom Rand der Modulbox beträgt halbe
Rahmenbreite. Für das Einfügen von Bildern etwa kann es sinnvoll sein,
diesen Rahmen wegzulassen. Dies geschieht mit der Sternchen-Variante
\Environment{posterrow*}.


\begin{minipage}{0.6\textwidth}
\begin{Example}
  \noindent\Macro{begin}\PParameter{\Environment{posterrow}}
    \OParameter{X,X}\par
  \noindent Teilt eine Modulzeile in 2 Spalten auf.
\end{Example}
\end{minipage}
\hfill
\begin{minipage}{0.3\textwidth}\centering
  \examplegraphic[width=\textwidth,page=2]{examples/module_split.pdf}
\end{minipage}

\begin{Declaration}
  \XMacro{begin}\PParameter{\Environment{postercol}}%
    \OParameter{rows}\\%
  \quad\dots\\
  \XMacro{end}\PParameter{postercol}\\
  \XMacro{begin}\PParameter{\Environment{postercol*}}%
    \OParameter{rows}\\%
  \quad\dots\\
  \XMacro{end}\PParameter{postercol*}
\end{Declaration}

Die mit \Environment{posterrow} angelegten Spalten können jeweils mit 
\Environment{postercol} mit Inhalt gefüllt oder in neue Unterzeilen aufgeteilt 
werden.

\begin{minipage}{0.6\textwidth}
\begin{Example}
  \noindent\Macro{begin}\PParameter{\Environment{postercol}}
    \OParameter{3cm,X}\par
  \noindent Teilt eine Modulspalte in 2 Unterzeilen auf.
\end{Example}
\end{minipage}
\hfill
\begin{minipage}{0.3\textwidth}\centering
  \examplegraphic[width=\textwidth,page=3]{examples/module_split.pdf}
\end{minipage}


\begin{Declaration}
  \XMacro{begin}\PParameter{\Environment{postersubrow}}\\%
  \quad\dots\\
  \XMacro{end}\PParameter{postersubrow}\\
  \XMacro{begin}\PParameter{\Environment{postersubrow*}}\\%
  \quad\dots\\
  \XMacro{end}\PParameter{postersubrow*}
\end{Declaration}

Die mit \Environment{postercol} angelegten Unterzeilen können jeweils mit 
\Environment{postercol} mit Inhalt gefüllt werden.

\begin{Example}

\end{Example}


\clearpage
\section{Aushänge}

\begin{lstlisting}
\documentclass[style=bulletin]{tubsposter}
\end{lstlisting}

% \section{Seitenlayout}

% margin
% bcor
% sender=top/bottom


% \section{Format}

% \section{Farben}

% \section{Befehle}

% \section{Optionen}
      % Kapitel: 'Plakate'

% \chapter{Broschüren}

Mit der Klasse \newdocumentclass{tubsleaflet} können dreigeteilte Broschüren
im Corporate Design erstellt werden.
Die Vorlage basiert auf der Klasse \texttt{leaflet}.
Alle allgemeinen Informationen und Möglichkeiten sind daher in der
entsprechenden Dokumentation\cite{cls:leaflet} nachzuschlagen.

Die Ausgabe erfolgt zweiseitig auf DINA4-Größe.
Andere Formate können gewählt werden, sind aber nicht voll unterstützt.

\begin{important}
  Als Papierformat sollte nicht \Option{a4paper} gewählt werden,
  da so falsche Größen geladen werden. Die korrekte Option lautet
  \Option{langpaper} (\gls{glos:dinlang}),
  ist aber optional, da sie automatisch vorgeladen wird.
\end{important}


\begin{figure}[!ht]
  \examplegraphic[width=\textwidth]{examples/leaflet.pdf}
  \caption{Einfacher Flyer im CD (erste Seite)}
\end{figure}

\paragraph{Darstellungsstil}\hfill

\begin{Declaration}
  \KOption{style}\PName{<simple/none>}
\end{Declaration}

Der Darstellungsstil \PName{simple} setzt standardmäßig das Frontseiten- und das
Rückseiten-Logo.
Im Darstellungsstil \PName{none} werden standardmäßig keine Elemente gesetzt.

Hintergrundelemente können manuell mittels \Macro{AddToBackground} gesetzt
werden.

\paragraph{Seiten}\hfill

Die einzelnen Seiten können mit der Umgebung \Environment{gausspage} erstellt
werden. Diese erlaubt das Erstellen von Seiten im Gaußraster.
Da sie ein universelles Element von \tubslatex ist, ist ihr ein spezielles
\href{chap:gausspage}{Kapitel} gewidmet.     % Kapitel: 'Broschüren'

% \chapter{Gauß-Layout}\label{chap:gausspage}

Dieses Kapitel beschreibt das allgemeine Layoutsystem zur Erstellung von
Inhalten und Hintergründen im Gaußraster des Corporate Designs.

\begin{Declaration}
  \Macro{showtubslogo}\OParameter{Position}
\end{Declaration}

Bewirkt Darstellung des TU-Siegelbandlogos im aktiven Layout.
Die Option \PName{Position} erlaubt die Angabe der Darstellungsseite
(links/rechts). Standardmäßig wird das Logo links bzw. innen dargestellt.

\begin{Declaration}
  \Macro{showlogo}\PParameter{Logo}
\end{Declaration}

Bewirkt Darstellung eines Individuellen Logos im aktuellen Layout.
\PName{Logo} kann dabei entweder einfacher Text oder auch ein 
mit \Macro{includegraphics} eingebundenes Bild sein.

Die Positionierung wird automatisch an die Positionierung des Siegelbandlogos
angepasst. Wird dies links bzw. innen platziert, so steht das individuelle
Logo rechts bzw. außen.

\begin{Declaration}
  \Macro{showtopline}
\end{Declaration}

Bewirkt Darstellung einer Trennlinie zwischen Absender und Kommunikationsbereich
im aktuellen Layout.

\begin{Declaration}
  \Macro{showdesignhelper}
\end{Declaration}

\begin{Declaration}
  \Macro{bgelement}\OParameter{Darstellung}\PParameter{Höhe}
\end{Declaration}

Erstellt ein Hintergrundelement im Gaußraster mit angegebener \PName{Höhe}.
Der Parameter \PName{Darstellung} kann die folgenden Einstellungen verarbeiten:

\begin{Declaration}
  \KOption{bgcolor}\PName{Farbe}\\
  \KOption{bgimage}\PName{Bild-Datei}\\
  \KOption{imagefit}\PName{Darstellungsoption}
\end{Declaration}

Mit \OptionValue{bgcolor}{Farbe} wird das Hintergrundelement mit der angegebenen
Farbe gefüllt.

Die Option \OptionValue{bgimage}{Bild-Datei} erlaubt dagegen die Darstellung
eines Hintergrundbildes im Element.
Da der Darstellungsbereich fest vorgegeben ist, muss das eingebundene Bild
in diesen Bereich eingepasst werden. Dies geschieht automatisch, die Art
der Einpassung lässt sich aber mit der Option \Option{imagefit} kontrollieren.
Sie erlaubt folgende Einstellungen:


\begin{desctable}
\entry{\PValue{cropped}}{%
  Automatisches Abschneiden. Dies ist die Standardeinstellung.
}
\entry{\PValue{cropx}}{%
  Abschnitt horizontal.
}
\entry{\PValue{cropy}}{%
  Abschnitt vertikal.
}
\entry{\PValue{scaled}}{%
  Horizontale \emph{und} vertikale Skalierung.
}
\end{desctable}
   % Kapitel: 'Gauß-Layout'

\part{Basis-Pakete}

Dieser Abschnitt befasst sich mit den Paketen, die von den jeweiligen
\tubslatex-Anwenderklassen normalerweise automatisch geladen werden.
Diese Pakete stellen einige grundlegende Funktionalitäten und Darstellungen
für das \acs{CD} der TU Braunschweig dar. Dazu gehören zum einen das Paket
\newpackage{tubscolors}, das alle Farben, die im \acs{CD} vorgesehen sind,
definiert. Des weiteren gibt es das Paket \newpackage{tubslogo}, welches
zur Darstellung des Siegelbandlogos dient. Abschließend ist noch das Paket
\newpackage{nexus} zu erwähnen, welches die Hausschrift \emph{Nexus}
lädt.

All diese Pakete sind so konzipiert, dass sie auch unabhängig von den
Klassen und Paketen aus \tubslatex eingesetzt werden können, auch wenn dies
im Allgemeinen nicht erforderlich sein sollte.

\chapter{Farben}\label{chap:tubscolors}

\newcommand{\classoptionitem}[1][ ]{
  \item[\mdseries{\ttfamily%
    \textbackslash usepackage%
    {[{\color{tuRed}#1}]}%
    \{tubslogo\}}]\hfill\\
}

Die Farbdefinitionen in \tubslatex werden vom Paket \newpackage{tubscolors}
zur Verfügung gestellt. Die folgende Beschreibung bezieht sich
auf den Funktionsumfang des Paketes. Das Paket ist in allen verfügbaren
Klassen bereits korrekt eingebunden.

\newcommand{\rainbow}[2][\relax]{{\noindent\sffamily\footnotesize%
\ifx#1\relax\colorlet{fglbg}{black}\else\colorlet{fglbg}{#1}\fi
\colorbox{#2100}{\hbox to 0.188\textwidth{%
  \color{fglbg}\vphantom{Fg}#2{}100\hfill}}% 
\colorbox{#280}{\hbox to 0.188\textwidth{%
  \color{fglbg}\vphantom{Fg}#2{}80\hfill}}% 
\colorbox{#260}{\hbox to 0.188\textwidth{\vphantom{Fg}#2{}60\hfill}}% 
\colorbox{#240}{\hbox to 0.188\textwidth{\vphantom{Fg}#2{}40\hfill}}% 
\colorbox{#220}{\hbox to 0.188\textwidth{\vphantom{Fg}#2{}20\hfill}}\\% 
}}

\section{Verfügbare Farben}

Der Farbklang der TU-Braunschweig ist in eine Primär- und einen
Sekundärfarbbereich aufgeteilt.

Die Primärfarben bilden dabei Rot, Schwarz und Weiß, sowie in
20-Prozent-Schritten abgestufte Grautöne. Die Primärfarben dienen
vor allem zur Auszeichnung von Hintergrund, Textfarbe und dem TU-Logo.
Zur individuellen Gestaltung von Dokumenten ist der Sekundärfarbbereich
vorgesehen.

Die Sekundärfarben setzen sich aus 12 weiteren aufeinander abgestimmten
Farben zusammen, die 
in 4 Farbklänge (Gelb-Orange, Grün, Blau und Violett) mit je 3 Basisfarben
aufgeteilt sind.
Alle Sekundärfarben können in 20-Prozent-Schritten aufgehellt werden.

Die Namen über die die einzelnen Farben angesprochen werden können, sind in den
Beispielfeldern angegeben.

\subsection{Primärfarben}

{\sffamily\footnotesize%
\colorbox{tuRed}{\hbox to 0.188\textwidth{%
  \vphantom{Fg}tuRed\hfill}}%
\colorbox{tuBlack}{\hbox to 0.188\textwidth{%
  \color{white}\vphantom{Fg}tuBlack\hfill}}%
\fcolorbox{tuBlack}{tuWhite}{\hbox to 0.188\textwidth{%
  \vphantom{Fg}tuWhite\hfill}}\\%
}

\rainbow[tuWhite]{tuGray}

Zur Vereinfachung sind noch die Farben \lstinline{tuGray} und
\lstinline{tuLightGrey} definiert, die den Farben \lstinline{tuGray60} und
\lstinline{tuGray20} entsprechen.

Alle Graytöne sind darüber hinaus auch in britischer Schreibweise nutzbar
(\lstinline{tuGrey}).

\paragraph{Hinweis:}
\lstinline{tuRed} ist nicht zu verwechseln mit \lstinline{tuRed100} aus dem
Sekundärfarbbereich. Es handelt sich dabei um eine komplett andere Farbe.

\pagebreak
\subsection{Sekundärfarben}

  tuOrange\ldots\\
  \colorshow{Orange}{Light}
  \colorshow{Orange}{Medium}
  \colorshow{Orange}{Dark}\\[-1ex]
  tuBlue\ldots\\
  \colorshow{Blue}{Light}
  \colorshow{Blue}{Medium}
  \colorshow{Blue}{Dark}\\[-1ex]
  tuGreen\ldots\\
  \colorshow{Green}{Light}
  \colorshow{Green}{Medium}
  \colorshow{Green}{Dark}\\[-1ex]
  tuViolet\ldots\\
  \colorshow{Violet}{Light}
  \colorshow{Violet}{Medium}
  \colorshow{Violet}{Dark}
%   \caption{Im CD definierte Farben und deren Benennung (Auszug)}


Zusätzlich kann als Paketoption ein Farbklang ausgewählt werden, dessen Farben
dann über die Werte  \\\lstinline{tuSecondaryLight},
\lstinline{tuSecondaryMedium}, \lstinline{tuSecondaryDark}, sowie die
entsprechenden Prozentualwert \\(\lstinline{tuSecondaryLight20},
\lstinline{tuSecondaryLight40}, \ldots) angesprochen werden können.
Dies erlaubt eine flexible Verwendung der 4 Sekundärfarbklänge.

In folgendem Beispiel wurde \lstinline{blau} als Farbklang ausgewählt:

\colorshow{Secondary}{Light}
\colorshow{Secondary}{Medium}
\colorshow{Secondary}{Dark}\\[-1ex]

\paragraph{Hinweise:}
Die Farben des gelb-orange-Farbklangs können entsprechend der anderen
Farbmodelle auch noch einheitlich über die Alternativnamen
\lstinline{tuOrangeLight}, \lstinline{tuOrange},
\lstinline{tuOrangeDark}, sowie de entsprechenden Prozentwerte aufgerufen
werden.

Außerdem können jeweils die mittleren Farbwerte der Farbklänge auch über den
Zusatz \lstinline{Medium} angesprochen werden (statt \lstinline{tuGreen100} auch
\lstinline{tuGreenMedium100}).

Bei allen 100-Prozent-Farben (außer \lstinline{tuRed}) kann die
Zahl weggelassen werden (statt \lstinline{tuGreenLight100} auch 
\lstinline{tuGreenLight}).

\subsection{Farbmodelle}

In den Paketoptionen kann zwischen 3 Farbmodellen gewählt werden.
Das Standardmodell stellt die CD-konformen RGB-Farbwerte zur Verfügung.

Für die Ausgabe in CMYK-Farben steht die Option \lstinline!cmyk! zur Verfügung.

Darüber hinaus gibt es noch ein RGB-Farbschema, das für die Ausgabe auf
Beamern optimiert ist. Es kann über die Option \lstinline!rgbbeamer! geladen
werden.
  % Kapitel: 'Farben'

\chapter{Das Siegelbandlogo}

Die verschiedenen Farbvarianten des Siegelbandlogos in \tubslatex werden
vom Paket \lstinline{tubslogo} zur Verfügung gestellt.
Die folgende Beschreibung bezieht sich auf den Funktionsumfang des Paketes.
Das Paket ist in allen verfügbaren Klassen bereits korrekt eingebunden.

Zur Verfügung gestellt werden eine Variante im rgb-Farbmodell, eine im
cmyk-Farbmodell, eine mit der Sonderfarbe HKS 15 oder (für Ausnahmefälle)
eine monochromen Variante.
Die Auwahl erfolgt dabei über Paketoptionen
(siehe Abschnitt \ref{options:color}).

\begin{figure}[!ht]
\begin{minipage}{0.5\textwidth}
  \centering
  \includegraphics[width=\tubslogoBaseWidth]{TUBraunschweig-rgb}
  \subcaption{Logo in Standardgröße (RGB)}
\end{minipage}
\begin{minipage}{0.5\textwidth}
  \centering
  \includegraphics[width=\tubslogoBaseWidth]{TUBraunschweig_4C}
  \subcaption{Logo in Standardgröße (CMYK)}
\end{minipage}

\begin{minipage}{0.5\textwidth}
  \centering
  \includegraphics[width=\tubslogoBaseWidth]{TUBraunschweig_SC}
  \subcaption{Logo in Standardgröße (HKS 15)}
\end{minipage}
\begin{minipage}{0.5\textwidth}
  \centering
  \includegraphics[width=\tubslogoBaseWidth]{TUBraunschweig_B}
  \subcaption{Logo in Standardgröße (schwarz)}
\end{minipage}
\caption{Verfügbare Farbversionen der Logos}
\end{figure}

\begin{Declaration}
  \Macro{tubslogo}\OParameter{Skalierung}\\
  \Macro{tubslogoAbs}\Parameter{Breite}
\end{Declaration}


Es werden 2 Befehle zum Einfügen des Logos in Dokumente bereit gestellt.
Mit dem Befehl \Macro{tubslogo} wird ein Logo in
standardkonformer Größe der verwendeten Papiergröße entsprechend skaliert
eingefügt. Die verwendete Papiergröße kann dabei als Paketoption übergeben
werden und ist standardmäßig auf \lstinline{a4paper} voreingestellt
(siehe Abschnitt \ref{options:papersize}).
Als optionales Argument kann auch ein frei gewählter Skalierungsfaktor
vorgegeben werden (siehe Abschnitt \ref{cmd:tubslogo}).

Mit \Macro{tubslogoAbs} kann über den Parameter \PName{Breite} dagegen
zusätzlich eine individuell gewählte absolute Breite angegeben werden.


\section{Paket-Optionen}

\subsection{Farbmodell}\label{options:color}

\begin{description}
  \classoptionitem[rgb]
    Verwendung der rgb-Druckfarben. Dies ist die Standardeinstellung.
  \classoptionitem[cmyk]
    Verwendung der cmyk-Farben
  \classoptionitem[hks]
    Verwendung der Spezialfarbe HKS 15.
  \classoptionitem[mono]
    Verwendung einer schwarz-weiß-Version des Logos.
\end{description}

\subsection{Papiergröße}\label{options:papersize}

\begin{description}
  \classoptionitem[a6paper]
    Skaliert das Logo passend für A6-Dokumente auf 60\% der Standardgröße.
  \classoptionitem[a5paper]
    Skaliert das Logo passend für A5-Dokumente auf 70\% der Standardgröße.
  \classoptionitem[a4paper]
    Skaliert das Logo passend für A4-Dokumente auf Standardgröße.
  \classoptionitem[a3paper]
    Skaliert das Logo passend für A3-Dokumente auf 140\% der Standardgröße.
  \classoptionitem[a2paper]
    Skaliert das Logo passend für A2-Dokumente auf 200\% der Standardgröße.
  \classoptionitem[a1paper]
    Skaliert das Logo passend für A1-Dokumente auf 280\% der Standardgröße.
  \classoptionitem[a0paper]
    Skaliert das Logo passend für A0-Dokumente auf 400\% der Standardgröße.
\end{description}


\clearpage
\section{Befehle}

\begin{description}\label{cmd:tubslogo}
  \item[\color{tuRed}\mdseries\ttfamily \textbackslash tubslogo%
    {[\textcolor{tuGreenDark}{\sffamily\itshape $<$Breite rel.$>$}]}]
    Erzeugt ein Logo, dass entsprechend des durch Angabe der Papiergröße
    gewählten Skalierungsfaktors skaliert wird.

    Wird das optionale Argument {\sffamily\itshape $<$Breite rel.$>$}
    verwendet, so kann ein freier Skalierungsfaktor ausgehend von der
    Standardgröße für DINA4-Dokumente ($70$mm$\times 26$mm) übergeben werden.

    \lstinline!\tubslogo[0.6]! erzeugt beispielsweise ein auf 60\% der 
    Standardgröße für DINA4-Dokumente skaliertes Logo.

  \item[\color{tuRed}\mdseries\ttfamily \textbackslash tubslogoAbs%
    \{\textcolor{tuGreenDark}{Breite abs.}\}]
    Dieser Befehl erlaubt die Erzeugung eines Logos mit absolut vorgegebener
    Breite.\footnote{Generell ist die Verwendung von standardkonformen Größen
    einer individuellen Skalierung vorzuziehen!}

    \lstinline!\tubslogoAbs{5cm}! erzeugt beispielsweise ein 5cm breites Logo.
\end{description}


\section{Längen}
  \begin{description}
    \item[\mdseries\ttfamily \textbackslash tubslogoBaseWidth]\hfill\\
      Standardbreite des Logos (70mm).
    \item[\mdseries\ttfamily \textbackslash tubslogoBaseHeight]\hfill\\
      Standardhöhe des Logos (26mm).
    \item[\mdseries\ttfamily \textbackslash tubslogoWidth]\hfill\\
      Breite des automatisch skalierten Logos. Wird automatisch berechnet und
      kann zur dynamischen Positionierung verwendet werden.
    \item[\mdseries\ttfamily \textbackslash tubslogoHeight]\hfill\\
      Höhe des automatisch skalierten Logos. Siehe \lstinline{tubslogoWidth}.
  \end{description}
    % Kapitel: 'Logo'

%%
%% Copyright (C) 2011-2015 by Enrico Jörns
%%

\input docstrip.tex
\keepsilent

\preamble

This is a generated file.

Copyright (C) 2011-2015 by Enrico Jörns <e.joerns@tu-bs.de>

This file may be distributed and/or modified under the
conditions of the LaTeX Project Public License, either
version 1.2 of this license or (at your option) any later
version. The latest version of this license is in:

  http://www.latex-project.org/lppl.txt

and version 1.2 or later is part of all distributions of
LaTeX version 1999/12/01 or later.

\endpreamble

% --- Prevent overwrite questions ---
\askforoverwritefalse

% --- File generation ---

\usedir{tex/latex/nexus}

\generate{\file{nexus.sty}{%
  \from{tubsvers.inc}{default,extra}
  \from{nexus.dtx}{package}
}}

\generate{\file{LY1NexusProSans.fd}{%
  \from{nexus.dtx}{fdsans}
}}

\generate{\file{LY1NexusProSans-lnum.fd}{%
  \from{nexus.dtx}{fdsanslnum}
}}

\generate{\file{LY1NexusProSerif.fd}{%
  \from{nexus.dtx}{fdserif}
}}

\generate{\file{LY1NexusProSerif-lnum.fd}{%
  \from{nexus.dtx}{fdseriflnum}
}}

\endbatchfile
   % Kapitel: 'Nexus'

\appendix

% \chapter{Anhang}

\def\foo{
\section{Kompatibilität zu anderen Paketen}

\begin{description}
  \item[fontenc]
    Das fontenc-Paket wird von allen Vorlagen automatisch im korrekten Encoding
    geladen uns sollte nicht manuell nachgeladen werden.
  \item[xcolor]
    Das xcolor-Paket wird von allen Vorlagen bereits eingebunden.
    % todo: table & other options?
\end{description}
}

\chapter{Längendefinitionen}

\begin{tabularx}{\textwidth}{lX}
  \Length{tubssenderheight}  &
    Höhe des Absenderbereichs\\
  \Length{tubscommunicationheight} &
    Höhe des Kommunikationsbereichs\\
  \Length{tubsborderwidth}  &
    Breite des Rahmens\\
%   \Length{tubssegmentheight} &
%     Höhe des aktuellen Segments\\
  \Length{tubscolumnsep} &
    Abstand von Spalten im Spaltenlayout\\
  \Length{tubsgaussbaseheight}  &
    Basishöhe des kleinsten Gauß-Segements\\
  \Length{tubscolumnwidth}  &
    Breite einer Spalte im Spaltenlayout.
\end{tabularx}


%Glossar ausgeben
\printglossary[style=altlist,title=Glossar]
%Abkürzungen ausgeben
\deftranslation[to=German]{Acronyms}{Abkürzungsverzeichnis}
\printglossary[type=\acronymtype,style=long]

\chapter{lco-Dateien}
\section{Institut}
\lstinputlisting{../../letter/doc/examples/musterinstitut.lco}
\section{Mitarbeiter}
\lstinputlisting{../../letter/doc/examples/mustermann.lco}


\listoffigures

\bibliographystyle{alpha}
\bibliography{lit}

\end{document}
